% Méthodes numériques pour physiciens et ingénieurs
% Woafo Paul -- transcription LaTeX du document scanné
\documentclass[11pt,a4paper,oneside]{report}

\usepackage[T1]{fontenc}
\usepackage[utf8]{inputenc}
\usepackage[french]{babel}
\usepackage{lmodern}
\usepackage{microtype}
\usepackage[a4paper,margin=2.5cm,headheight=14pt]{geometry}
\usepackage{amsmath,amssymb,mathtools}
\usepackage{array,multirow,booktabs}
\usepackage{arydshln}
\usepackage{enumitem}
\usepackage{titlesec}
\usepackage{xcolor}
\usepackage{graphicx}
\usepackage{tikz}
\usetikzlibrary{arrows.meta,calc,patterns,decorations.pathreplacing}
\usepackage{pgfplots}
\pgfplotsset{compat=1.18}
\usepackage{algorithm}
\usepackage{algpseudocode}
\usepackage[most]{tcolorbox}
\usepackage{fancyhdr}
\usepackage[titles]{tocloft}
\setlength{\cftchapnumwidth}{2.6em}
\setlength{\cftsecindent}{1.5em}
\setlength{\cftsecnumwidth}{3.4em}
\setlength{\cftsubsecindent}{4.9em}
\setlength{\cftsubsecnumwidth}{4.4em}
\usepackage[hidelinks]{hyperref}

% ---------------------------------------------------------------
% Numérotation à la manière du document original : I-1-, I-2-1-, ...
% ---------------------------------------------------------------
\renewcommand{\thechapter}{\Roman{chapter}}
\renewcommand{\thesection}{\thechapter-\arabic{section}}
\renewcommand{\thesubsection}{\thesection-\arabic{subsection}}
\renewcommand{\thesubsubsection}{\thesubsection-\arabic{subsubsection}}
\setcounter{secnumdepth}{3}
\setcounter{tocdepth}{2}

\titleformat{\chapter}[display]
  {\normalfont\bfseries\centering}
  {\LARGE Chapitre \thechapter}{0.6em}{\Huge}
\titlespacing*{\chapter}{0pt}{-20pt}{30pt}
\titleformat{\section}{\normalfont\Large\bfseries}{\thesection-}{0.8em}{}
\titleformat{\subsection}{\normalfont\large\bfseries}{\thesubsection-}{0.6em}{}
\titleformat{\subsubsection}{\normalfont\normalsize\bfseries}{\thesubsubsection-}{0.5em}{}
\titleformat{\paragraph}[runin]{\normalfont\normalsize\bfseries\itshape}{}{0pt}{}[.]

\numberwithin{equation}{chapter}
\renewcommand{\theequation}{\thechapter.\arabic{equation}}

% ---------------------------------------------------------------
% Algorithmes en français
% ---------------------------------------------------------------
\floatname{algorithm}{Algorithme}
\renewcommand{\listalgorithmname}{Liste des algorithmes}
\algrenewcommand\algorithmicrequire{\textbf{Données :}}
\algrenewcommand\algorithmicensure{\textbf{Résultat :}}
\algrenewcommand\algorithmicfor{\textbf{Pour}}
\algrenewcommand\algorithmicdo{\textbf{faire}}
\algrenewcommand\algorithmicwhile{\textbf{Tant que}}
\algrenewcommand\algorithmicif{\textbf{Si}}
\algrenewcommand\algorithmicthen{\textbf{alors}}
\algrenewcommand\algorithmicelse{\textbf{Sinon}}
\algrenewcommand\algorithmicrepeat{\textbf{Répéter}}
\algrenewcommand\algorithmicuntil{\textbf{Jusqu'à}}
\algrenewcommand\algorithmicreturn{\textbf{Retourner}}
\algrenewcommand\algorithmicend{\textbf{Fin}}
\algrenewtext{EndFor}{\textbf{Fin pour}}
\algrenewtext{EndWhile}{\textbf{Fin tant que}}
\algrenewtext{EndIf}{\textbf{Fin si}}
\algnewcommand\algorithmicprint{\textbf{Écrire}}
\algnewcommand\Print{\State\algorithmicprint{} }
\algnewcommand\algorithmicread{\textbf{Lire}}
\algnewcommand\Read{\State\algorithmicread{} }
\algnewcommand{\Goto}[1]{\State\textbf{Aller à} #1}
\algnewcommand{\Stop}{\State\textbf{Arrêt}}
\algrenewcommand\algorithmiccomment[1]{\hfill\textit{\{#1\}}}

% Environnement pour les algorithmes non flottants
\newenvironment{algo}[1][]
  {\begin{tcolorbox}[enhanced,breakable,colback=gray!4,colframe=gray!60,
     boxrule=0.4pt,arc=1mm,left=6pt,right=6pt,top=4pt,bottom=4pt,
     title={#1},fonttitle=\bfseries\small,coltitle=black,colbacktitle=gray!15]
   \setlength{\lineskip}{6pt}\setlength{\lineskiplimit}{3pt}%
   \begin{algorithmic}[0]}
  {\end{algorithmic}\end{tcolorbox}}

% ---------------------------------------------------------------
% Encadrés « Exercice(s) » et « Remarque »
% ---------------------------------------------------------------
\newtcolorbox{exercice}[1][Exercice]{enhanced,breakable,colback=blue!3,
  colframe=blue!40!black,boxrule=0.4pt,arc=1mm,
  title={#1},fonttitle=\bfseries,left=6pt,right=6pt}

\newcommand{\NB}{\textbf{NB :}~}

% ---------------------------------------------------------------
% Macros mathématiques
% ---------------------------------------------------------------
\newcommand{\dd}{\mathrm{d}}
\newcommand{\R}{\mathbb{R}}
\newcommand{\C}{\mathbb{C}}
\renewcommand{\Re}{\operatorname{Re}}
\renewcommand{\Im}{\operatorname{Im}}
\DeclareMathOperator{\signe}{signe}
\newcommand{\abs}[1]{\left|#1\right|}

\setlength{\parskip}{0.4em}
\setlength{\parindent}{1.5em}
\allowdisplaybreaks

\pagestyle{fancy}
\fancyhf{}
\fancyhead[L]{\small\itshape\leftmark}
\fancyfoot[C]{\thepage}
\renewcommand{\headrulewidth}{0.4pt}
\renewcommand{\chaptermark}[1]{\markboth{Chapitre \thechapter~: #1}{}}
\fancypagestyle{plain}{\fancyhf{}\fancyfoot[C]{\thepage}\renewcommand{\headrulewidth}{0pt}}

\hypersetup{
  pdftitle={Méthodes numériques pour physiciens et ingénieurs},
  pdfauthor={Woafo Paul},
  pdfsubject={Méthodes numériques},
  pdflang={fr},
  hypertexnames=false
}

\begin{document}

\pagenumbering{roman}
% ---------------------------------------------------------------
% Page de titre
% ---------------------------------------------------------------
\hypersetup{pageanchor=false}
\begin{titlepage}
  \centering
  \vspace*{2cm}
  {\rule{\linewidth}{1.5pt}\par}
  \vspace{0.8cm}
  {\fontsize{34}{40}\selectfont\bfseries MÉTHODES NUMÉRIQUES\par}
  \vspace{0.8cm}
  {\rule{\linewidth}{1.5pt}\par}
  \vspace{1.5cm}
  {\LARGE\bfseries Pour physiciens\par}
  \vspace{0.3cm}
  {\Large et\par}
  \vspace{0.3cm}
  {\LARGE\bfseries Ingénieurs\par}
  \vfill
  {\Large\bfseries WOAFO Paul\par}
  \vspace{0.3cm}
  {\large Professeur, Université de Yaoundé I\par}
  \vspace{1.2cm}
  {\Large\bfseries Cameroun\par}
  \vspace{2cm}
\end{titlepage}
\hypersetup{pageanchor=true}

\chapter*{Préface}
\addcontentsline{toc}{chapter}{Préface}
\markboth{Préface}{}

Le but de ce document est de permettre aux scientifiques en général, aux
étudiants physiciens et ingénieurs en particulier, d'élaborer des méthodes de
calcul numérique utilisables par l'ordinateur pour résoudre les problèmes de
recherche et de développement.

Il s'agit d'une compilation des notes de cours dispensés par le Professeur
WOAFO Paul de 1994 à 2016. C'est en effet en 1994 que cet enseignement a été
introduit dans les programmes de formation des étudiants de Maîtrise de
Physique au Cameroun (à l'Université de Yaoundé, seule université
camerounaise en ce moment-là). Plus tard, un enseignement de méthodes
numériques a été introduit en année de Licence de Physique. Et les travaux
pratiques étaient effectués au Centre de Calcul (aujourd'hui Centre
Universitaire des Technologies de l'Information).

Année après année, l'enseignement s'est enrichi et a contribué à la formation
des étudiants de Maîtrise, Master et Doctorat camerounais qui pour la plupart
étaient engagés dans des travaux théoriques pouvant avoir une forte composante
numérique. Certains de ces étudiants occupent actuellement des positions
académiques au Cameroun et à l'étranger. Certains d'entre eux se sont inspirés
des éléments de ce cours pour former d'autres étudiants dans les universités
camerounaises. L'importance historique de cet enseignement dans la production
scientifique des physiciens camerounais est donc incommensurable.

Ces notes des cours ont bénéficié des contributions de nombreux
collaborateurs, notamment, Professeur NANA NBENDJO Blaise (Faculté des
Sciences, Université de Yaoundé~I), Dr.~NOUBISSIE Samuel (IUT Fotso-Victor de
Bandjoun, Université de Dschang) et Dr.~FAUTSO Gaétan (École Normale
Supérieure de Bambili, Université de Bamenda).

Il faut noter que l'implémentation des méthodes numériques a beaucoup évolué
ces dernières années avec la mise sur le marché des logiciels comme MATLAB
contenant des codes pour les différentes méthodes numériques. Mais, la plupart
des problèmes de recherche et de développement sont complexes et exigent une
expertise en matière de méthodes numériques (schémas numériques et
algorithmes). De plus, pour utiliser les logiciels existant sur le marché, on
doit être capable de comprendre les différentes étapes afin de savoir comment
et où fournir des informations aux logiciels. Il ne faut pas aussi oublier le
prix d'achat des logiciels.

Il nous a donc paru important de mettre ces notes de cours à la disposition
des étudiants. Cette importance est renforcée par la croissance rapide des
effectifs des étudiants ; ce qui a un effet négatif sur les conditions de
formation.

Le document présente les méthodes numériques de base qui sont dispensées en
année de Licence et de Master de Physique et des études en sciences de
l'ingénieur. Il comporte les méthodes numériques pour :
\begin{itemize}[label=--]
  \item le traitement numérique des mesures expérimentales,
  \item le calcul des intégrales et des dérivées,
  \item la résolution des systèmes d'équations algébriques linéaires et non
        linéaires,
  \item la résolution des équations différentielles et
        intégro-différentielles,
  \item la résolution des équations aux dérivées partielles.
\end{itemize}

Les méthodes de simulation numérique des équations issues des problèmes
d'advection-diffusion, d'écoulement des fluides (équations de Navier-Stokes)
et la dynamique des structures atomiques ne sont pas incluses dans ce
document. Elles pourront éventuellement être présentées en séminaires ou de
manière détaillée aux étudiants de 2\textsuperscript{e} année de Master ou de
Doctorat. Il en est de même des approches Méthodes des Éléments Finis et des
Volumes Finis.

Chaque partie du cours fournit d'une part les développements théoriques des
méthodes numériques, et d'autre part des algorithmes correspondants (prêts à
être traduits dans un langage de programmation tels que le Pascal, le Fortran
ou le C, que les étudiants sont supposés bien maîtriser). Une extension de
l'utilisation de certaines méthodes numériques expliquées dans ce document est
destinée à la programmation des microcontrôleurs qui sont susceptibles de
convertir les instructions du programme en signaux électriques directement
utilisables pour des applications en mécatronique ou dans les systèmes
automatiques ou en électronique pour la génération des signaux.

La synthèse présentée dans ce document a bénéficié des supports didactiques
déjà disponibles pour l'enseignement des méthodes numériques, dont notamment
les livres suivants : S.~Vandergraft, \og Introduction to Numerical
Computations\fg{}, Academic Press (New York, 1978) ; R.L.~Burden,
J.D.~Faires, A.C.~Reynolds, \og Numerical Analysis\fg{}, 2\textsuperscript{nd}
edition, Weber and Schmidt (Boston, 1981) et D.~Euvrard, \og Résolution
numérique des équations aux dérivées partielles\fg{}, Masson (Paris, 1987).

\begin{flushright}
  \textit{Woafo Paul, Professeur}
\end{flushright}


\tableofcontents
\clearpage
\pagenumbering{arabic}

\chapter{Systèmes d'équations algébriques linéaires}

\section{Introduction}

Le traitement des problèmes scientifiques conduit très souvent à des systèmes
d'équations algébriques linéaires à plusieurs inconnues. Lorsque le nombre
d'équations n'est pas élevé (par exemple moins de 3), on peut trouver les
solutions du système d'équations par des méthodes analytiques connues
(éliminations, substitution). Mais lorsque ce nombre devient élevé ou lorsque
les coefficients des équations sont variables, il est utile de développer des
codes de calcul numériques qui aideront l'ordinateur à calculer les solutions.

L'objet de ce chapitre est de présenter les schémas numériques pour la
résolution des équations algébriques linéaires : les méthodes directes
(élimination de Gauss et de Jordan, racine carrée de Cholesky, décomposition
triangulaire, fragmentation) et les méthodes itératives (de Jacobi et de
Gauss-Seidel).

\section{Méthodes directes pour les systèmes d'équations linéaires}

\subsection{Méthodes d'élimination de Gauss et de Jordan}

\subsubsection{Principe de la méthode}

Soit le système algébrique de la forme :
\[
  AX = b
\]
où $A = (a_{ij})$, $X = (x_i)$ et $b = (b_i)$.

Ce système d'équations est de la forme
\[
  \left\{
  \begin{aligned}
    a_{11}x_1 + a_{12}x_2 + \dots + a_{1n}x_n &= b_1\\
    a_{21}x_1 + a_{22}x_2 + \dots + a_{2n}x_n &= b_2\\
    &\;\;\vdots\\
    a_{n1}x_1 + a_{n2}x_2 + \dots + a_{nn}x_n &= b_n
  \end{aligned}
  \right.
\]

La méthode d'élimination de Gauss est basée sur le principe suivant. Les
inconnues sont éliminées successivement en exprimant à partir d'une des
équations une inconnue en fonction des autres puis en substituant dans le
reste des équations (voir classes du secondaire). On part ainsi d'un système de
$n$ équations à $n$ inconnues à un système de $n-1$ équations à $n-1$
inconnues. Le processus se poursuit jusqu'à l'obtention d'une équation à une
inconnue. On trouve alors cette dernière inconnue et par un processus inverse,
les autres inconnues sont calculées. À chaque étape, l'équation qui permet
d'exprimer une inconnue en fonction des autres est appelée \emph{équation
pivotale}. Pour éviter les erreurs d'arrondis, l'équation pivotale doit être
celle qui possède le plus grand coefficient de l'inconnue à éliminer.

La procédure de cette méthode peut alors se formuler de la manière suivante :

Pour une ligne n\textsuperscript{o}~$i$, notée $L_i$, on fait l'opération
\[
  L_i \leftarrow L_i - \frac{a_{i1}}{a_{11}}L_1, \qquad i = 2 \text{ jusqu'à } n.
\]

Pour $i$ allant de $2$ à $n$, on obtient alors un nouveau système de la forme
\[
  \left\{
  \begin{aligned}
    a_{11}x_1 + a_{12}x_2 + \dots + a_{1n}x_n &= b_1\\
    0 + a'_{22}x_2 + \dots + a'_{2n}x_n &= b'_2\\
    0 + a'_{32}x_2 + \dots + a'_{3n}x_n &= b'_3\\
    &\;\;\vdots\\
    0 + a'_{n2}x_2 + \dots + a'_{nn}x_n &= b'_n
  \end{aligned}
  \right.
\]
Les coefficients $a'_{ij}$ et $b'_i$ sont donnés par
\[
  a'_{ij} = a_{ij} - \frac{a_{i1}}{a_{11}}a_{1j}, \qquad
  b'_i = b_i - \frac{a_{i1}}{a_{11}}b_1 .
\]
On continue le processus en considérant le système constitué des lignes $2$ à
$n$ et ainsi de suite jusqu'au système constitué des lignes $n-1$ et $n$.

\subsubsection{Algorithme}

Ainsi pour une étape $k$, c'est-à-dire en considérant le système constitué des
lignes $k$ à $n$, l'algorithme suivant est utilisé :

\begin{algo}
  \For{$l$ allant de $k+1$ à $n$}
    \State $b_l \leftarrow b_l - \dfrac{a_{lk}}{a_{kk}}\,b_k$
    \For{$j$ allant de $k+1$ à $n$}
      \State $a_{lj} \leftarrow a_{lj} - \dfrac{a_{lk}}{a_{kk}}\,a_{kj}$
    \EndFor
  \EndFor
\end{algo}

On peut ainsi en déduire l'algorithme de triangularisation du système
d'équations en faisant varier $k$ de $1$ à $n-1$. Le système triangulaire se
présente ainsi
\[
  \left\{
  \begin{array}{l@{\;}l@{\;}l@{\;}c@{\;}l}
    a_{11}x_1 + & a_{12}x_2 + & \dots\dots\dots\dots\dots\dots & + a_{1n}x_n &= b_1\\
    0 + & a_{22}x_2 + & \dots\dots\dots\dots\dots\dots & + a_{2n}x_n &= b_2\\
    0 + & 0 & + a_{33}x_3 \dots\dots\dots\dots & + a_{3n}x_n &= b_3\\
    \vdots & & & & \\
    0 + & 0 + & \dots\dots\dots\dots\dots\dots + 0 & + a_{nn}x_n &= b_n
  \end{array}
  \right.
\]

\NB Il faut noter que les coefficients $a_{ij}$ et $b_k$ du système
triangulaire ne sont pas ceux du départ.

Le système algébrique triangulaire peut se résoudre de la manière suivante :

La dernière équation donne
\[
  x_n = \frac{b_n}{a_{nn}} .
\]
Pour l'inconnue $x_k$, on a
\[
  x_k = \bigl(b_k - a_{k,k+1}x_{k+1} - a_{k,k+2}x_{k+2} - \dots - a_{k,n}x_n\bigr)\big/ a_{kk}
\]
ou sous forme condensée
\[
  x_k = \Bigl(b_k - \sum_{l=k+1}^{n} a_{k,l}\,x_l\Bigr)\Big/ a_{kk} .
\]

\begin{exercice}[Exercices]
\begin{enumerate}
  \item Écrire l'algorithme de résolution du système algébrique triangulaire.
  \item Écrire l'algorithme qui permet d'interchanger les lignes des systèmes
        algébriques dû à la règle de pivot maximal.
  \item Écrire l'algorithme complet de la méthode de Gauss (triangularisation
        et résolution du système triangulaire) en prenant soin de faire jouer
        la règle du pivot maximal.
\end{enumerate}
\end{exercice}

\subsubsection{Méthode de Jordan}

Elle est analogue à la méthode de Gauss. Mais ici, on met aussi à zéro les
termes qui se trouvent au-dessus de la diagonale principale. On obtient alors
un système diagonal.

\subsection{Méthodes par décomposition de la matrice}

\subsubsection{Méthode de la décomposition triangulaire}

\paragraph{a- Principe}

Soit le système
\[
  AX = b .
\]
On peut de manière unique décomposer la matrice $A$ en deux matrices
triangulaires dont l'une est triangulaire inférieure avec les éléments
quelconques sur la diagonale et l'autre triangulaire supérieure avec tous les
éléments égaux à 1 sur la diagonale.

Soit $L$ et $S$ ces deux matrices triangulaires, on a
\[
  A = LS .
\]
Le système d'équations devient
\[
  SX = L^{-1}b .
\]
Ce dernier système est un système triangulaire que l'on peut résoudre par le
schéma précédent.

\paragraph{Décomposition de la matrice}

On peut écrire $L$ et $S$ sous la forme
\[
  L = \begin{pmatrix}
    l_{11} & 0 & 0 & \cdots & 0\\
    l_{21} & l_{22} & 0 & \cdots & 0\\
    \vdots & & \ddots & & \vdots\\
    \vdots & & & \ddots & 0\\
    l_{n1} & l_{n2} & l_{n3} & \cdots & l_{nn}
  \end{pmatrix},
  \qquad
  S = \begin{pmatrix}
    1 & s_{12} & s_{13} & \cdots & s_{1n}\\
    0 & 1 & s_{23} & \cdots & s_{2n}\\
    \vdots & & 1 & & s_{3n}\\
    \vdots & & & \ddots & \vdots\\
    0 & 0 & 0 & \cdots & 1
  \end{pmatrix}.
\]
L'équation $LS = A$ entraîne par identification les relations suivantes
\[
\begin{array}{ll}
  \left\{\begin{aligned}
    a_{j1} &= l_{j1}\\
    a_{1j} &= l_{11}s_{1j}
  \end{aligned}\right.
  &
  \begin{aligned}
    &1 \le j \le n\\
    &2 \le j \le n
  \end{aligned}
  \\[1.2em]
  \left\{\begin{aligned}
    a_{j2} &= l_{j1}s_{12} + l_{j2}\\
    a_{2j} &= l_{21}s_{1j} + l_{22}s_{2j}
  \end{aligned}\right.
  &
  \begin{aligned}
    &2 \le j \le n\\
    &3 \le j \le n
  \end{aligned}
  \\[1.2em]
  \left\{\begin{aligned}
    a_{j3} &= l_{j1}s_{13} + l_{j2}s_{23} + l_{j3}\\
    a_{3j} &= l_{31}s_{1j} + l_{32}s_{2j} + l_{33}s_{3j}
  \end{aligned}\right.
  &
  \begin{aligned}
    &3 \le j \le n\\
    &4 \le j \le n
  \end{aligned}
  \\
  \quad\vdots & \\
  \left\{\begin{aligned}
    a_{j,n-1} &= l_{j1}s_{1,n-1} + l_{j2}s_{2,n-1} + \dots + l_{j,n-2}s_{n-2,n-1} + l_{j,n-1}\\
    a_{n-1,j} &= l_{n-1,1}s_{1j} + l_{n-1,2}s_{2j} + \dots + l_{n-1,n-1}s_{n-1,j}
  \end{aligned}\right.
  &
  \begin{aligned}
    &n-1 \le j \le n\\
    &j = n
  \end{aligned}
  \\[1.2em]
  \;a_{nn} = l_{n1}s_{1n} + l_{n2}s_{2n} + \dots + l_{nn} &
\end{array}
\]

Avec ces équations on obtient les éléments de la première colonne de $L$ puis
ceux de la première ligne de $S$, ensuite ceux de la 2\textsuperscript{e}
colonne de $L$ et ceux de la 2\textsuperscript{e} ligne de $S$ et ainsi de
suite. On obtient alors les éléments suivants qui permettent de calculer les
éléments $l_{ij}$ et $s_{ij}$. On a :
\[
\begin{aligned}
  &\left\{\begin{aligned}
    l_{j1} &= a_{j1}\\
    s_{1j} &= \frac{a_{1j}}{l_{11}}
  \end{aligned}\right.
  \qquad
  \left\{\begin{aligned}
    l_{j2} &= a_{j2} - l_{j1}s_{12}\\
    s_{2j} &= \frac{1}{l_{22}}\bigl(a_{2j} - l_{21}s_{1j}\bigr)
  \end{aligned}\right.
  \\[0.6em]
  &\left\{\begin{aligned}
    l_{j3} &= a_{j3} - l_{j1}s_{13} - l_{j2}s_{23}\\
    s_{3j} &= \frac{1}{l_{33}}\bigl(a_{3j} - l_{31}s_{1j} - l_{32}s_{2j}\bigr)
  \end{aligned}\right.
  \\
  &\quad\vdots\\
  &\left\{\begin{aligned}
    l_{j,n-1} &= a_{j,n-1} - \sum_{k=1}^{n-2} l_{jk}\,s_{k,n-1}\\
    s_{n-1,j} &= \frac{1}{l_{n-1,n-1}}\Bigl[a_{n-1,j} - \sum_{k=1}^{n-2} l_{n-1,k}\,s_{kj}\Bigr]
  \end{aligned}\right.
  \\[0.6em]
  &\; l_{nn} = a_{nn} - \sum_{k=1}^{n-1} l_{nk}\,s_{kn}
\end{aligned}
\]

\begin{exercice}
\begin{enumerate}
  \item Écrire l'algorithme permettant de calculer les $l_{ij}$ et $s_{ij}$.
  \item Trouver dans la littérature l'algorithme d'inversion d'une matrice et
        de calcul des éléments du produit $L^{-1}b$.
\end{enumerate}
\end{exercice}

\subsubsection{Méthode de la racine carrée de Cholesky}

Une matrice $A$ symétrique peut être mise sous la forme d'un produit de deux
matrices $M$ et $M'$ dont l'une est la transposée de l'autre. On a ainsi une
seule matrice triangulaire supérieure $M$ à calculer. Ceci simplifie les
calculs et diminue l'encombrement mémoire de la machine. En utilisant le
procédé précédent on établit que les éléments de $M$ sont définis par :
\[
\begin{aligned}
  m_{11} &= \sqrt{a_{11}}\\
  m_{1j} &= \frac{a_{1j}}{m_{11}}\\
  m_{22} &= \sqrt{a_{22} - \sum_{k=1}^{1} m_{k2}^2} = \sqrt{a_{22} - m_{12}^2}\\
  m_{2j} &= \frac{a_{2j} - \sum_{k=1}^{1} m_{k2}m_{kj}}{m_{22}}
          = \frac{a_{2j} - m_{12}m_{1j}}{m_{22}}\\
  m_{pp} &= \sqrt{a_{pp} - \sum_{k=1}^{p-1} m_{kp}^2}
          && \text{pour } p \text{ allant de } 2 \text{ à } n\\
  m_{pj} &= \frac{a_{pj} - \sum_{k=1}^{p-1} m_{kp}m_{kj}}{m_{pp}}
          && \text{pour } p \text{ allant de } 2 \text{ à } n
             \text{ et } j \text{ allant de } p+1 \text{ à } n\\
  m_{nn} &= \sqrt{a_{nn} - m_{1n}^2 - m_{2n}^2 - \dots - m_{n-1,n}^2}
\end{aligned}
\]

\begin{exercice}
Écrire l'algorithme de calcul des coefficients $m_{pp}$ et $m_{pj}$.
\end{exercice}

\subsubsection{Fragmentation des systèmes volumineux}

Quand $n$ est très grand, il n'y a pas assez de mémoire pour stocker les
$n(n+1)$ coefficients du système et mettre en place le programme de calcul. On
fragmente alors le système d'équations en décomposant la matrice du système en
sous-matrices.

Par exemple, on mettra le système
\[
  AX = b
\]
sous la forme
\[
  {\renewcommand{\arraystretch}{1.5}
  \left(\begin{array}{c:c}
    A_1 & A_2\\ \hdashline
    A_3 & A_4
  \end{array}\right)
  \left(\begin{array}{c}
    X_1\\ \hdashline
    X_2
  \end{array}\right)
  =
  \left(\begin{array}{c}
    B_1\\ \hdashline
    B_2
  \end{array}\right)}.
\]
Les matrices $A_1$ et $A_2$ ne sont pas nécessairement des matrices carrées et
la décomposition par des droites en pointillés est telle que les
multiplications matricielles suivantes sont possibles :
\[
  \left\{
  \begin{aligned}
    A_1X_1 + A_2X_2 &= B_1\\
    A_3X_1 + A_4X_2 &= B_2
  \end{aligned}
  \right.
\]
Si $A_1$ est régulière, alors on a
\[
  X_1 = A_1^{-1}\bigl[B_1 - A_2X_2\bigr].
\]
En substituant dans la deuxième équation, on obtient
\[
  \bigl[A_3A_1^{-1}A_2 - A_4\bigr]X_2 = \bigl[A_3A_1^{-1}B_1 - B_2\bigr].
\]
On peut ainsi calculer les éléments de $X_2$ et substituer dans l'expression
de $X_1$ pour avoir les éléments de $X_1$. On obtient donc ainsi toutes les
composantes de $X$.

Cette méthode peut être généralisée en décomposant la matrice en $9$, $16$,
$25$, \dots{} sous-matrices.

\subsubsection{Systèmes à coefficients complexes}

Soit le système matriciel
\[
  AZ = b
\]
avec
\[
  A = M + iN, \qquad b = c + id, \qquad Z = X + iY .
\]
Le système algébrique peut alors se décomposer en deux sous-systèmes de la
forme
\[
  \left\{
  \begin{aligned}
    MX - NY &= c\\
    NX + MY &= d
  \end{aligned}
  \right.
\]
On peut alors réunir ces deux sous-systèmes en un système unique, ou utiliser
l'approche du paragraphe précédent.

\section{Méthodes itératives pour les systèmes d'équations linéaires}

\subsection{Formules itératives}

Les méthodes directes telles que la méthode d'élimination de Gauss sont
appropriées pour les systèmes de taille moyenne. Pour les systèmes de grande
taille, il est conseillé d'utiliser les méthodes itératives qui sont plus
rapides, plus faciles à programmer et moins encombrantes. Cependant, ces
méthodes itératives sont basées sur le calcul des suites infinies qu'il faudra
tronquer à partir d'un certain rang. Pour le choix des valeurs initiales des
itérations, on fixe généralement les inconnues $x_i = 1$ ou $x_i = 0$ pour
$i = 1$ jusqu'à $n$.

Le principe des méthodes itératives consiste à exprimer l'inconnue $x_i$ de la
ligne $i$ en fonction des autres inconnues ; c'est-à-dire
\[
  x_i = \frac{b_i - \displaystyle\sum_{\substack{j=1\\ j\neq i}}^{n} a_{ij}x_j}{a_{ii}} .
\]
Ceci suggère que les valeurs connues $x_1, x_2, x_3, \dots, x_{i-1}, x_{i+1},
\dots, x_n$ soient substituées pour trouver $x_i$. Au départ, il faut donc
donner des valeurs initiales aux inconnues $x_i$. Ces valeurs initiales sont
généralement fixées à $x_i = 1$ ou $x_i = 0$ pour $i$ allant de $1$ à $n$.

Il existe deux variantes de méthode itérative : celle de \textbf{Jacobi} et
celle de \textbf{Gauss-Seidel}.

La formule itérative de \textbf{Jacobi} est :
\[
  x_i^{(k+1)} = \frac{b_i - \displaystyle\sum_{\substack{j=1\\ j\neq i}}^{n} a_{ij}x_j^{(k)}}{a_{ii}}
\]
où $x_i^{(k)}$ sont les valeurs des inconnues après $k$ itérations.

Dans la formule de \textbf{Gauss-Seidel}, les valeurs $x_j^{(k+1)}$ pour
$j \le i-1$ sont utilisées aussitôt qu'elles sont déterminées pour calculer
$x_i^{(k+1)}$. La formule itérative de \textbf{Gauss-Seidel} se met alors sous
la forme :
\[
  x_i^{(k+1)} = \frac{b_i - \displaystyle\sum_{j=1}^{i-1} a_{ij}x_j^{(k+1)}
                 - \displaystyle\sum_{j=i+1}^{n} a_{ij}x_j^{(k)}}{a_{ii}} .
\]
Notons que les deux formules sont importantes car il existe des systèmes
d'équations pour lesquels la méthode de Jacobi converge tandis que la méthode
de Gauss-Seidel ne converge pas. Dans d'autres cas c'est la méthode de
Gauss-Seidel qui converge et celle de Jacobi ne converge pas.

\subsection{Conditions d'arrêt}

Pour la condition d'arrêt de l'itération, il y a plusieurs méthodes. Elles
consistent à arrêter les itérations lorsque toutes les quantités
\[
  d_i = \abs{x_i^{(k+1)} - x_i^{(k)}} < \varepsilon
\]
ou aussi
\[
  d_i = \abs{\frac{x_i^{(k+1)} - x_i^{(k)}}{x_i^{(k+1)}}} < \varepsilon .
\]
Cependant, le meilleur test de convergence est d'utiliser la condition
\[
  d = \frac{\displaystyle\sum_{i=1}^{n}\abs{x_i^{(k+1)} - x_i^{(k)}}}
           {\displaystyle\sum_{i=1}^{n}\abs{x_i^{(k+1)}}} < \varepsilon .
\]

\begin{exercice}
Écrire les algorithmes des méthodes de Jacobi et de Gauss-Seidel avec
imposition de la condition d'arrêt.
\end{exercice}

\chapter{Équations algébriques non linéaires}

\section{Introduction}

Les équations algébriques non linéaires se présentent sous la forme
$f(x) = 0$ ou $f_i(x) = 0$ avec $i$ allant de $1$ à $n$. La variable $x$ peut
être réelle ou complexe. La résolution numérique de ces équations fait appel à
plusieurs méthodes qui seront expliquées dans ce chapitre. Il s'agit notamment
de la méthode de Newton-Raphson, la méthode de dichotomie et de la méthode de
Bairstow. Ce chapitre expliquera aussi comment effectuer les opérations sur
les polynômes ainsi que la présentation des méthodes permettant de trouver
plusieurs racines simultanément. Le cas des systèmes d'équations non linéaires
sera aussi présenté.

\section{Méthodes de Newton-Raphson et dichotomie}

La formule de Newton-Raphson a été donnée par Raphson en 1698. Mais Newton
avait suggéré une méthode plus voisine quelques mois auparavant. Il se
pourrait aussi que cette formule ait été utilisée par Héron en l'an 100 avant
Jésus-Christ pour trouver la racine carrée d'un nombre. Cette méthode est de
la classe des méthodes itératives par approximations successives qui se
traduisent par répétition systématique d'une opération déterminée, ce qui
finit par donner de proche en proche la solution exacte.

\subsection{Formule de la méthode de Newton-Raphson}

On commence généralement par trouver une solution grossière ou approchée
$x_0$, dite approximation d'ordre zéro. Dans le cas d'un problème de physique,
une telle solution $x_0$ peut être suggérée par la signification du problème
ou par une courbe. Généralement, on localise un intervalle $[a,b]$ tel que
$x_0$ appartienne à cet intervalle. Si nous désignons la solution exacte par
$\bar{x} = x_0 + h$, alors, la formule de Taylor donne
\[
  f(x_0 + h) \approx f(x_0) + h f'(x_0) + O(h^2)
\]
or
\[
  f(x_0 + h) = 0 \;\longrightarrow\; f(x_0) + h f'(x_0) \approx 0
  \;\longrightarrow\; h \approx -\frac{f(x_0)}{f'(x_0)}
\]
d'où
\[
  \bar{x} = x_0 - \frac{f(x_0)}{f'(x_0)} .
\]
En posant $x_1 = x_0 - \dfrac{f(x_0)}{f'(x_0)}$, on obtient une première
approximation de la solution (l'indice désignant l'ordre de l'approximation).
Procédons de même avec $x_1$, on obtient la seconde approximation,
\[
  x_2 = x_1 - \frac{f(x_1)}{f'(x_1)} .
\]
À l'ordre $n+1$, on obtient
\[
  x_{n+1} = x_n - \frac{f(x_n)}{f'(x_n)} .
\]
C'est une formule résultant d'une erreur d'ordre $O(h^2)$. En négligeant les
termes d'ordre supérieur, on a remplacé la courbe représentative de $f(x)$ par
une tangente en $x = x_0$ ; la méthode consiste géométriquement à tracer de
proche en proche les tangentes à la courbe et à rechercher leur point
d'intersection avec l'axe des abscisses. En effet, on a l'équation
\[
  f(\bar{x}) = y = f(x_0) + h f'(x_0) = f(x_0) + (\bar{x} - x_0) f'(x_0)
\]
qui est l'équation de la tangente au point $\bigl(x_0, f(x_0)\bigr)$.

Cette interprétation graphique fait que la méthode de Newton est souvent
appelée \emph{méthode de la tangente}.

\paragraph{Remarques} Lorsque $f(x_0)f''(x_0) > 0$, la suite
$x_{n+1} = x_n - \dfrac{f(x_n)}{f'(x_n)}$ est monotone, bornée et converge vers
$\bar{x}$. Ce résultat donne une condition suffisante pour que la méthode de
Newton converge vers $\bar{x}$. Mais cette condition n'est pas nécessaire.
Dans le cas où il n'y a pas convergence, on change la valeur de $x_0$.

\subsection{Condition d'arrêt}

La méthode de Newton converge vers la solution exacte $\bar{x}$, mais on ne
connaît pas explicitement à quel moment une certaine précision est atteinte.
Pour cela, on utilise généralement le résultat suivant :

Soit $\bar{x}$ une racine de l'équation $f(x) = 0$ et soit $x_n$ une valeur
approchée de $\bar{x}$. Si sur $[a,b]$ contenant $\bar{x}$ et $x_n$, on a
\[
  \abs{f'(x)} \ge m > 0, \quad \text{alors} \quad
  \abs{x_n - \bar{x}} \le \frac{\abs{f(x_n)}}{m} .
\]
Dans la pratique, on pourra prendre $m$ comme
$m = \min\bigl(\abs{f'(a)}, \abs{f'(b)}\bigr)$.

Si $E$ est la précision exigée, on s'arrêtera lorsque
$\dfrac{\abs{f(x_n)}}{m} < E$.

\begin{exercice}
Écrire l'algorithme de la méthode de Newton-Raphson en y intégrant la
condition d'arrêt.
\end{exercice}

\subsection{Autres formules : formule du troisième et du premier ordre}

La formule de Newton est l'une des méthodes les plus utilisées parce qu'elle
est relativement simple et d'une convergence très suffisante dans la plupart
des cas. Il faut toutefois que $x_0$ soit assez près de la racine pour que la
dérivée $f'$ ne change pas de signe entre $x_0$ et $x$. Si en effet $f'$
s'annule sans que $f$ s'annule aussi, la tangente à la courbe va couper l'axe
des $x$ à l'infini. On risque alors d'osciller indéfiniment autour de la
solution sans jamais l'atteindre.

Il est possible de modifier la formule de Newton et d'écrire
\[
  x_{n+1} = x_n - p\,\frac{f(x_n)}{f'(x_n)}
\]
où $p$ est un nombre ou une fonction. On peut aussi utiliser la formule
\[
  x_{n+1} = x_n - \frac{g\,f(x_n)}{g\,f'(x_n) + h\,f(x_n)} .
\]
En développant la série de Taylor jusqu'à l'ordre 2, on obtient
\[
  f(x_0 + h) = f(x_0) + h f'(x_0) + \frac{h^2}{2} f''(x_0) + O(h^3) .
\]
En remplaçant $h$ par $h = \bar{x} - x_0 \approx -\dfrac{f}{f'}$ (d'après la
formule de Newton), on obtient
\[
  x_{n+1} = x_n - \frac{f(x_n)}{f'(x_n)}
          - \frac{1}{2}\,\frac{f''(x_n)}{f'(x_n)}
            \left(\frac{f(x_n)}{f'(x_n)}\right)^{2} .
\]
C'est la formule du troisième ordre.

Soit l'équation $f(x) = 0$. De cette équation, on peut tirer $x$ en fonction
de $x$ en écrivant $x = g(x)$, où $g$ est défini à partir de $f(x)$. Dans ce
cas on peut écrire $x_{n+1} = g(x_n)$. C'est la formule d'itération du premier
ordre.

\subsection{Formule de Newton pour les racines complexes}

Soit l'équation
\[
  f(z) = 0
\]
avec $z = x + iy$. On a
\[
  f(z) \approx f(z_0) + (z - z_0) f'(z_0)
  \;\Longrightarrow\; z - z_0 = -\frac{f(z_0)}{f'(z_0)} .
\]
Pour le calcul numérique, il faut séparer les parties réelles et imaginaires.
On aura donc
\[
  f(z_0) = G(x_0, y_0) + iH(x_0, y_0), \qquad
  \frac{\dd f}{\dd z_0} = f'(z_0) = G'(x_0, y_0) + iH'(x_0, y_0)
\]
et
\[
  \left\{
  \begin{aligned}
    x_{n+1} &= x_n - \left.\frac{GG' + HH'}{G'^2 + H'^2}\right|_{z_n}\\[0.4em]
    y_{n+1} &= y_n - \left.\frac{HG' - GH'}{G'^2 + H'^2}\right|_{z_n}
  \end{aligned}
  \right.
\]

\subsection{Méthode de dichotomie ou de bissection}

Supposons que $f(a) < 0 < f(b)$ et $f$ est strictement croissante sur $[a,b]$.
La solution $x_0$ de l'équation $f(x) = 0$ appartient à l'intervalle $[a,b]$.
Le milieu de l'intervalle est $c = (a+b)/2$. Si $f(a)$ et $f(c)$ ont des signes
opposés, alors $x_0$ appartient à $[a,c]$. Si par contre $f(b)$ et $f(c)$ sont
de signes opposés, alors $x_0$ appartient à $[c,b]$. Dans l'un ou l'autre
cas, la taille de l'intervalle est divisée par 2. Une nouvelle bissection peut
être faite sur le nouvel intervalle contenant la solution et ainsi de suite.
Ainsi, après $n$ étapes, la taille de l'intervalle devient $(b-a)/2^n$.

\begin{exercice}
Écrire l'algorithme de la méthode de bissection.
\end{exercice}

\section{Calcul de toutes les racines d'une équation polynomiale}

\subsection{Opérations sur les polynômes}

\subsubsection{Schéma de calcul numérique d'un polynôme}

Pour déterminer les valeurs numériques de
\[
  P_n(x) = a_0x^n + a_1x^{n-1} + \dots + a_{n-1}x + a_n ,
\]
on peut calculer chaque terme et faire la somme. Mais, ce procédé n'est pas en
général court. Il exige $2n-1$ multiplications. On peut utiliser le schéma dit
de \textbf{Horner}. On calcule les nombres $b_k$ définis par :
\[
  \begin{aligned}
    b_0 &= a_0,\\
    b_k &= a_k + b_{k-1}x \qquad (k \text{ allant de } 1 \text{ à } n).
  \end{aligned}
\]
$b_n$ est alors la valeur numérique cherchée.

L'algorithme utilisé se présente ainsi :

\begin{algo}
  \State \textbf{Début}
  \State \quad $x \leftarrow$ \textit{valeur de} $x$
  \State \quad $b_0 \leftarrow a_0$
  \State \quad \textbf{Pour} $k$ allant de $1$ à $n$ \textbf{faire}
  \State \quad\quad $b_k \leftarrow a_k + b_{k-1}x$
  \State \quad \textbf{Fin pour}
  \State \quad \textbf{Écrire} $b_n = P_n(x)$
  \State \textbf{Fin}
\end{algo}

La première méthode de calcul exige $(2n-1)$ multiplications alors que le
schéma de Horner n'en demande que $n$ multiplications et le nombre d'additions
est le même dans les deux cas. De plus le schéma de Horner a l'avantage de
n'introduire qu'un type d'opérations :
\[
  b_k = a_k + b_{k-1}x .
\]
La programmation est alors plus aisée.

Pour un polynôme à coefficients complexes,
\[
  \begin{aligned}
    &P_n(z) = c_0z^n + c_1z^{n-1} + \dots + c_{n-1}z + c_n,\\
    &z = x + iy \quad\text{et}\quad c_p = a_p + ib_p,\\
    &P_n(z) = G_n(x,y) + iH_n(x,y).
  \end{aligned}
\]
Le schéma de Horner peut également être utilisé. On remplace tout simplement
les quantités $b_k$ par les quantités
\[
  P_k = G_k + iH_k
\]
avec
\[
  \begin{aligned}
    G_0 &= a_0, & H_0 &= b_0,\\
    G_k &= a_k + xG_{k-1} - yH_{k-1}, &
    H_k &= b_k + yG_{k-1} + xH_{k-1}.
  \end{aligned}
\]

\subsubsection{Division d'un polynôme par \texorpdfstring{$x-r$}{x-r}}

En divisant le polynôme
\[
  P_n(x) = a_0x^n + a_1x^{n-1} + \dots + a_{n-1}x + a_n
\]
par
\[
  C(x) = x - r ,
\]
on obtient un quotient $Q(x)$ et un reste $R$ liés entre eux par la relation
\[
  P(x) = (x - r)\,Q(x) + R .
\]
$Q(x)$ est un polynôme de degré $n-1$ et $R$ est une constante. Ils sont de la
forme
\[
  \begin{aligned}
    Q(x) &= b_0x^{n-1} + b_1x^{n-2} + \dots + b_{n-2}x + b_{n-1}\\
    R &= b_n
  \end{aligned}
\]
où $b_k = a_k + r\,b_{k-1}$ $(k = 1 \text{ à } n)$, avec $b_0 = a_0$.

Si $b_n = 0$, alors le polynôme est divisible par $x - r$ ; c'est-à-dire que
$r$ est une racine de $P(x)$.

\paragraph{Remarques}
\begin{itemize}
  \item[$*$] Ces formules sont semblables pour des grandeurs complexes. La
        seule différence apparaît dans le calcul effectif qui exige un
        dédoublement de colonne comme ci-dessus.
  \item[$**$] Quand les coefficients $a_k$ sont réels et que
        $r = \alpha + i\beta$, on peut diviser $P(x)$ par le facteur
        quadratique $(x - r)(x - r^*)$ où $r^*$ est le complexe conjugué.
\end{itemize}

\subsubsection{Division d'un polynôme par un facteur quadratique}

Soit le polynôme $P_n(x)$ que l'on désire diviser par le facteur quadratique
\[
  Q_2(x) = x^2 - Sx + P .
\]
On a :
\[
  P_n(x) = Q_2(x)\,Q_{n-2}(x) + R_1(x) .
\]
$Q_{n-2}(x)$ et $R_1(x)$ sont respectivement les polynômes de degré $n-2$ et
$1$ que l'on peut écrire sous la forme
\[
  \begin{aligned}
    Q_{n-2}(x) &= b_0x^{n-2} + b_1x^{n-3} + \dots + b_{n-3}x + b_{n-2}\\
    R_1(x) &= b_{n-1}(x - S) + b_n
  \end{aligned}
\]
avec
\[
  \begin{aligned}
    b_0 &= a_0, \qquad b_1 = a_1 + Sb_0,\\
    b_2 &= a_2 + Sb_1 - Pb_0,\\
    &\;\;\vdots\\
    b_k &= a_k + Sb_{k-1} - Pb_{k-2} \qquad (k = 2 \text{ à } n).
  \end{aligned}
\]
Si $S$ et $P$ sont respectivement la somme et le produit de deux racines du
polynôme $P_n(x)$, alors le reste $R_1$ est nul ; c'est-à-dire $b_{n-1} = 0$
et $b_n = 0$.

\subsection{Calcul d'une racine réelle et passage à l'équation de degré
\texorpdfstring{$n-1$}{n-1}}

L'idée la plus simple est de calculer la suite des valeurs
$b_k = a_k + b_{k-1}x_0$ qui conduisent à la valeur du polynôme et comme on
devrait avoir $b_n = 0$, on obtiendra
\[
  x = -\frac{a_n}{b_{n-1}} .
\]
On recommence le calcul des quantités $b_k$ avec cette nouvelle valeur de $x$
jusqu'à obtenir la précision choisie. À ce moment les coefficients $b_0$,
$b_1$, \dots, $b_{n-1}$ donnent le polynôme restant de degré $n-1$. On peut
alors continuer le processus de recherche des racines sur le polynôme de
degré $n-1$.
\[
  P_n(x) = a_0x^n + a_1x^{n-1} + \dots + a_{n-1}x + a_n
\]
L'inconvénient de ce processus est qu'il n'est pas toujours convergent. Cette
méthode appelée méthode de \textbf{Lin} fut développée en 1941.

Elle peut également s'appliquer sur le plan complexe. Pour cela, on part d'une
valeur arbitraire $z_0 = x_0 + iy_0$ et on calcule les coefficients de Horner
complexes jusqu'à $G_{n-1}$ et $H_{n-1}$. On obtient alors des nouvelles
valeurs $z_1 = x_1 + iy_1$ en écrivant que $G_n = 0$ et $H_n = 0$. On a en
effet
\[
  \left\{
  \begin{aligned}
    a_n + x_1G_{n-1} - y_1H_{n-1} &= 0\\
    b_n + x_1H_{n-1} + y_1G_{n-1} &= 0
  \end{aligned}
  \right.
  \quad\Longrightarrow\quad
  \left\{
  \begin{aligned}
    x_1 &= -\frac{a_nG_{n-1} + b_nH_{n-1}}{G_{n-1}^2 + H_{n-1}^2}\\[0.4em]
    y_1 &= \frac{a_nH_{n-1} - b_nG_{n-1}}{G_{n-1}^2 + H_{n-1}^2}
  \end{aligned}
  \right.
\]
On recommence le processus jusqu'à ce que
\[
  \abs{z^{(k+1)} - z^{(k)}} < \varepsilon .
\]

\begin{exercice}
Établir l'algorithme de ce schéma de Lin pour les solutions réelles.
\end{exercice}

\subsection{Formule d'itération de Newton sur l'axe réel}

Soit l'équation $P_n(x) = 0$, la dérivée de ce polynôme en $x$ est $P'_n(x)$.
Les valeurs numériques de $P_n(x)$ et de $P'_n(x)$ se calculent par le schéma
de Horner suivant :
\[
  \begin{aligned}
    b_0 &= a_0\\
    b_k &= a_k + xb_{k-1} \qquad (k = 1, \dots, n)
  \end{aligned}
\]
et
\[
  \begin{aligned}
    c_0 &= b_0\\
    c_k &= b_k + xc_{k-1} \qquad (k = 1, \dots, n-1)
  \end{aligned}
\]
On a
\[
  \left\{
  \begin{aligned}
    P_n(x) &= b_n\\
    P'_n(x) &= c_{n-1}
  \end{aligned}
  \right.
\]
La formule d'itération de Newton devient
\[
  x_{k+1} = f(x_k) \qquad\text{où}\qquad f(x) = x - \frac{b_n}{c_{n-1}} .
\]
On peut alors poursuivre le processus jusqu'à obtenir
\[
  \abs{x_{k+1} - x_k} < \varepsilon .
\]
Après avoir trouvé une racine $x^*$ de l'équation, on divise le polynôme par
$x - x^*$ et on continue le processus jusqu'à l'obtention de toutes les
racines.

\subsection{Calcul de deux racines et passage à l'équation de degré
\texorpdfstring{$n-2$}{n-2}}

Dans le cas de la division d'un polynôme par le facteur quadratique
\[
  G_2(x) = x^2 - Sx + P ,
\]
si $S$ et $P$ sont la somme et le produit de deux racines du polynôme
$P_n(x)$, alors on aura
\[
  b_n = 0 \quad\text{et}\quad b_{n-1} = 0
  \quad\Longrightarrow\quad
  P = \frac{a_n}{b_{n-2}} \quad\text{et}\quad
  S = \frac{Pb_{n-3} - a_{n-1}}{b_{n-2}} .
\]
Le processus peut alors continuer jusqu'à $b_n = b_{n-1} = 0$ à la précision
choisie. On obtient finalement un nouveau polynôme de degré $n-2$ pour lequel
on peut reprendre le calcul pour retrouver les deux racines suivantes et ainsi
de suite, on parvient à résoudre l'équation.

\subsection{Méthode de Bairstow}

Soit le polynôme :
\[
  P_n(x) = a_0x^n + a_1x^{n-1} + \dots + a_{n-1}x + a_n
\]
à diviser par $G_2(x) = x^2 - Sx + P$. On obtient
\[
  P_n(x) = G_2(x)\,Q_{n-2}(x) + R_1(x)
\]
où
\[
  \begin{aligned}
    Q_{n-2}(x) &= b_0x^{n-2} + b_1x^{n-3} + \dots + b_{n-3}x + b_{n-2}\\
    \text{et}\quad R_1(x) &= b_{n-1}(x - S) + b_n
  \end{aligned}
\]
avec
\[
  b_0 = a_0, \qquad b_1 = a_1 + Sb_0, \qquad
  b_k = a_k + Sb_{k-1} - Pb_{k-2} \quad (k = 2, \dots, n).
\]
Si l'on a $b_n = 0$ et $b_{n-1} = 0$, alors $S$ et $P$ sont la somme et le
produit de deux racines du polynôme et les grandeurs $b_k$ sont les
coefficients du polynôme de degré $n-2$. La méthode de Bairstow consiste à
partir des valeurs arbitraires de $S$ et $P$ et de calculer la suite des
coefficients $b_k$.

Soit
\[
  \begin{aligned}
    F(S,P) &= b_{n-1},\\
    G(S,P) &= b_n .
  \end{aligned}
\]
Si $F$ et $G$ ne sont pas nuls à la précision choisie, alors on améliore $S_0$
et $P_0$ ($S_0$ et $P_0$ étant les valeurs initiales de $S$ et de $P$) par la
formule d'itération de Newton à deux variables. Les deux fonctions de $S$ et
de $P$ sont $F$ et $G$ et les formules d'itération se mettent sous la forme
\[
  \begin{aligned}
    S_{k+1} &= S_k + \frac{\alpha}{\Delta}\\
    P_{k+1} &= P_k + \frac{\beta}{\Delta}
  \end{aligned}
\]
avec
\[
  \begin{aligned}
    \alpha &= F\frac{\partial G}{\partial P} - G\frac{\partial F}{\partial P}\\
    \beta &= G\frac{\partial F}{\partial S} - F\frac{\partial G}{\partial S}\\
    \Delta &= \frac{\partial G}{\partial S}\frac{\partial F}{\partial P}
            - \frac{\partial F}{\partial S}\frac{\partial G}{\partial P}
  \end{aligned}
\]
Pour obtenir les dérivées partielles par rapport à $S$, il suffit de dériver
les relations $b_k$ par rapport à $S$. Pour cela nous posons
\[
  \frac{\partial b_k}{\partial S} = c_{k-1} .
\]
En dérivant $b_k = a_k + Sb_{k-1} - Pb_{k-2}$ par rapport à $S$, on a
\[
  \frac{\partial b_k}{\partial S}
  = b_{k-1} + S\frac{\partial b_{k-1}}{\partial S} - P\frac{\partial b_{k-2}}{\partial S},
  \qquad\text{soit}\qquad
  c_{k-1} = b_{k-1} + Sc_{k-2} - Pc_{k-3} .
\]
On obtient alors la suite suivante
\[
  \begin{aligned}
    c_0 &= b_0\\
    c_1 &= b_1 + Sc_0\\
    c_k &= b_k + Sc_{k-1} - Pc_{k-2} \qquad (k = 2, \dots, n)
  \end{aligned}
\]
Posons
\[
  \frac{\partial b_k}{\partial P} = -c_{k-2} .
\]
On constate que l'on obtient une suite identique à celle définie ci-dessus.

Il vient finalement :
\[
  \begin{aligned}
    \alpha &= b_nc_{n-3} - b_{n-1}c_{n-2}\\
    \beta &= b_nc_{n-2} - b_{n-1}c_{n-1}\\
    \Delta &= c_{n-2}^2 - c_{n-1}c_{n-3}
  \end{aligned}
\]
Ces formules ont été établies par Bairstow en 1914 pour des besoins de
l'industrie d'aéronautique. Avec ces nouvelles valeurs de $S$ et de $P$, on
calcule à nouveau les coefficients $b_q$ et $c_q$ ainsi de suite jusqu'à ce
que $b_{n-1}$ et $b_n$ soient nuls à la précision choisie. Une fois que l'on a
terminé le calcul de $S$ et de $P$, on détermine les deux racines
correspondantes par :
\[
  x = \frac{1}{2}\Bigl[S \pm \sqrt{S^2 - 4P}\Bigr]
  \qquad\text{ou}\qquad
  x + iy = \frac{1}{2}\Bigl[S \pm i\sqrt{-S^2 + 4P}\Bigr] .
\]
Les coefficients restants $b_k$ sont ceux de l'équation de degré $n-2$. On peut
alors recommencer le processus et calculer par paire toutes les racines.
Toutefois, si $n$ est impair, il restera à la fin une équation de la forme :
\[
  b_0 + b_1x = 0 \quad\Longrightarrow\quad x = -\frac{b_0}{b_1} .
\]

\subsection{Calcul simultané de toutes les racines}

Pour calculer simultanément toutes les racines $x_1, x_2, \dots, x_n$ d'un
polynôme de degré $n$, il faut disposer de $n$ équations entre ces racines.
Malheureusement, c'est souvent difficile d'obtenir ce système d'équations
entre les racines. Mais pour les équations de degré inférieur à 5, on peut
trouver des relations : relations de Viète.

\subsubsection{Cas où toutes les racines sont réelles}

Si nous considérons l'équation du 3\textsuperscript{e} degré
\[
  x^3 + a_1x^2 + a_2x + a_3 = 0 .
\]
On a les relations de Viète suivantes entre les racines :
\[
  \left\{
  \begin{aligned}
    x_1 + x_2 + x_3 &= -a_1\\
    x_1x_2 + x_1x_3 + x_2x_3 &= a_2\\
    x_1x_2x_3 &= -a_3
  \end{aligned}
  \right.
  \quad\Longleftrightarrow\quad
  \left\{
  \begin{aligned}
    x_1 &= -a_1 - x_2 - x_3\\
    x_2 &= \frac{1}{x_1}\bigl(a_2 - x_1x_3 - x_2x_3\bigr)\\
    x_3 &= -\frac{a_3}{x_1x_2}
  \end{aligned}
  \right.
\]
On obtient ainsi un schéma itératif du 1\textsuperscript{er} ordre qui permet
par approximations successives d'avoir les trois racines de l'équation.
\[
  \left\{
  \begin{aligned}
    x_1^{(k+1)} &= -a_1 - x_2^{(k)} - x_3^{(k)}\\
    x_2^{(k+1)} &= \frac{1}{x_1^{(k)}}\Bigl(a_2 - x_1^{(k)}x_3^{(k)} - x_2^{(k)}x_3^{(k)}\Bigr)\\
    x_3^{(k+1)} &= -\frac{a_3}{x_1^{(k)}x_2^{(k)}}
  \end{aligned}
  \right.
\]

\subsubsection{Cas où il y a des racines complexes}

Si les coefficients de l'équation sont réels, les racines complexes se
présentent par paires conjuguées l'une de l'autre. En associant les racines
par couple et en faisant apparaître les sommes et les produits, on opère
uniquement sur les quantités réelles.

Pour une équation du 3\textsuperscript{e} degré :
$x^3 + a_1x^2 + a_2x + a_3 = 0$, si on pose $x_1 + x_2 = S$, $x_1x_2 = P$,
$x_3 = r$, les relations de Viète s'écrivent :
\[
  \left\{
  \begin{aligned}
    S + r + a_1 &= 0\\
    P + rS - a_2 &= 0\\
    rP + a_3 &= 0
  \end{aligned}
  \right.
\]
De ces relations, on tire :
\[
  \left\{
  \begin{aligned}
    S &= -r - a_1\\
    P &= a_2 - rS\\
    r &= -\frac{a_3}{P}
  \end{aligned}
  \right.
\]
On peut alors partir des valeurs quelconques de $S$, $P$ et $r$ pour faire le
processus itératif et déterminer ces grandeurs à la précision choisie.

\subsubsection{Extension aux équations de degré supérieur}

Pour une équation du 4\textsuperscript{e} ordre :
$x^4 + a_1x^3 + a_2x^2 + a_3x + a_4 = 0$, on a les relations suivantes :
\[
  \left\{
  \begin{aligned}
    S_1 + S_2 + a_1 &= 0\\
    S_1S_2 + P_1 + P_2 - a_2 &= 0\\
    S_1P_2 + S_2P_1 + a_3 &= 0\\
    P_1P_2 - a_4 &= 0
  \end{aligned}
  \right.
\]
avec
\[
  S_1 = x_1 + x_2, \quad P_1 = x_1x_2, \quad S_2 = x_3 + x_4
  \quad\text{et}\quad P_2 = x_3x_4 .
\]
On peut alors tirer le schéma itératif du 1\textsuperscript{er} ordre
suivant :
\[
  \left\{
  \begin{aligned}
    S_1 &= -a_1 - S_2\\
    P_1 &= a_2 - P_2 - S_1S_2\\
    S_2 &= -\frac{1}{P_1}\bigl(a_3 + S_1P_2\bigr)\\
    P_2 &= \frac{a_4}{P_1}
  \end{aligned}
  \right.
\]
\[
  \begin{aligned}
    x_1 &= \frac{S_1 - \sqrt{S_1^2 - 4P_1}}{2}\,; &
    x_2 &= \frac{S_1 + \sqrt{S_1^2 - 4P_1}}{2}\\[0.4em]
    x_3 &= \frac{S_2 - \sqrt{S_2^2 - 4P_2}}{2}\,; &
    x_4 &= \frac{S_2 + \sqrt{S_2^2 - 4P_2}}{2}
  \end{aligned}
\]
Pour une équation du 5\textsuperscript{e} degré on a :
\[
  \left\{
  \begin{aligned}
    S_1 &= -\bigl(a_1 + S_2 + r\bigr)\\
    P_1 &= a_2 - P_2 - S_1S_2 - r\bigl(S_1 + S_2\bigr)\\
    S_2 &= -\frac{1}{P_1}\Bigl[a_3 + S_1P_2 + r\bigl(P_1 + P_2 + S_1S_2\bigr)\Bigr]\\
    P_2 &= \frac{1}{P_1}\Bigl[a_4 - r\bigl(S_1P_2 + S_2P_1\bigr)\Bigr]\\
    r &= -\frac{a_5}{P_1P_2}
  \end{aligned}
  \right.
  \qquad\text{où } r = x_5 .
\]

\subsubsection{Quelques relations et règles importantes}

\paragraph{Relations de Viète}
Soit le polynôme de degré $n$ en $z$ défini par :
\[
  P_n(z) = a_0z^n + a_1z^{n-1} + \dots + a_{n-1}z + a_n
\]
où les $a_i$ et $z$ sont des réels ou complexes.

Soit $z_j$ les racines de ce polynôme, on a alors
\[
  P_n(z) = a_0(z - z_1)(z - z_2)\cdots(z - z_n) .
\]
En développant cette expression et après identification des coefficients on
obtient les relations suivantes :
\[
  \begin{aligned}
    a_1 &= -a_0\sum_i z_i\\
    a_2 &= a_0\sum_{i<j} z_iz_j\\
    a_3 &= -a_0\sum_{i<j<k} z_iz_jz_k\\
    &\;\;\vdots\\
    a_n &= (-1)^n a_0\,z_1z_2\cdots z_n
  \end{aligned}
\]

\paragraph{Localisation des racines dans le plan}

\subparagraph{-- Règle des signes de Descartes.}
Le nombre de racines réelles positives ne peut pas dépasser le nombre de
changements de signe que l'on observe en parcourant la suite des coefficients
du polynôme. Il ne peut différer de ce nombre que par un nombre pair. On a une
règle analogue pour les racines négatives. Il suffit de changer $x$ en $-x$.

\textbf{Exemples :} L'équation $x^4 - 8x^3 + 24x^2 - 32x + 15 = 0$ n'a pas de
racines réelles négatives. De même, l'équation
$x^4 - 2x^3 - 3x^2 + 8x - 4 = 0$ n'a pas plus de trois racines réelles
positives.

\subparagraph{-- Règle de Gua.}
Si le carré d'un coefficient intermédiaire est inférieur ou égal au produit
des coefficients voisins, alors il y a des racines imaginaires.

\textbf{Exemple :} L'équation $x^3 + 7x^2 + x + 8 = 0$ possède une racine
imaginaire car on a
\[
  1 \times 1 < 7 \times 8 .
\]

\subparagraph{-- Règle des lacunes.}
S'il manque un terme entre deux termes de même signe ou plusieurs termes entre
deux termes de signe quelconque, alors il y a des racines imaginaires.

\textbf{Exemple :} L'équation $x^3 + 6x + 20 = 0$ a le terme en $x^2$ absent
et les coefficients des termes $x^3$ et $x$ sont positifs. On peut donc
affirmer qu'il y a des racines complexes.

\subparagraph{-- Règle de Sturm.}
Soit $a$ un nombre réel positif, si l'équation $(x - a)P_n(x) = 0$ présente
$(2k+1)$ variations de signe de plus que l'équation $P_n(x) = 0$ alors il y a
au moins $2k$ racines complexes.

\textbf{Exemple :} Soit l'équation $x^3 + 7x^2 + x + 7 = 0$ et l'équation
$(x - 1)P_n(x) = 0$ égale à
\[
  x^4 + 6x^3 - 6x^2 + 6x - 7 = 0
\]
qui présente trois changements de signe. Il y a donc 2 racines complexes pour
l'équation $P_n(x) = 0$.

\section{Système d'équations non linéaires}

Soit un système d'équations non linéaires
\[
  f_j(x_i) = 0 \qquad
  \left\{
  \begin{aligned}
    i &= 1, \dots, n\\
    j &= 1, \dots, n
  \end{aligned}
  \right.
\]
On peut aussi élaborer des processus itératifs du 1\textsuperscript{er},
2\textsuperscript{e} et 3\textsuperscript{e} ordre pour ce type de système
d'équations. En particulier, pour le système de deux équations non linéaires
\[
  \left\{
  \begin{aligned}
    f_1(x,y) &= 0\\
    f_2(x,y) &= 0
  \end{aligned}
  \right. ,
\]
on peut trouver $x$ à partir de $f_1(x,y)$ et $y$ à partir de $f_2(x,y)$ et
utiliser le processus itératif du 1\textsuperscript{er} ordre suivant :
\[
  \left\{
  \begin{aligned}
    x_{k+1} &= F_1(x_k, y_k)\\
    y_{k+1} &= F_2(x_k, y_k)
  \end{aligned}
  \right.
\]
On peut également généraliser les formules itératives de Jacobi et de
Gauss-Seidel. De même, à partir de la méthode de Newton, le système de deux
équations non linéaires à 2 inconnues donne les formules itératives
suivantes :
\[
  \begin{aligned}
    x_{k+1} &= x_k + \frac{1}{\Delta}\bigl(f_2f_{1,y} - f_1f_{2,y}\bigr)_{(x_k,y_k)}
      \qquad\text{(quantité évaluée au point } (x_k,y_k)\text{)}\\
    y_{k+1} &= y_k + \frac{1}{\Delta}\bigl(f_1f_{2,x} - f_2f_{1,x}\bigr)_{(x_k,y_k)}\\
    \Delta &= f_{1,x}f_{2,y} - f_{1,y}f_{2,x}, \qquad
    f_{1,x} = \frac{\partial f_1}{\partial x}, \;\dots
  \end{aligned}
\]
Dans le cas d'un système de degré supérieur à 2, on élabore le schéma de
Newton de la manière suivante. Supposons que la solution d'ordre zéro soit le
vecteur $x_i^{(0)}$ et que la solution d'ordre 1 est $x_i^{(1)}$. On a
\[
  f_j\bigl(x^{(1)}\bigr) = f_j\bigl(x^{(0)}\bigr)
  + \sum_{i=1}^{n}\bigl(x_i^{(1)} - x_i^{(0)}\bigr)
    \left.\frac{\partial f_j}{\partial x_i}\right|_{x^{(0)}}
  \qquad (\text{évalué en } x^{(0)}) .
\]
Or $x_i^{(1)}$ étant solution de l'équation, on a $f_j\bigl(x^{(1)}\bigr) = 0$ ;
d'où le système matriciel
\[
  \sum_{i=1}^{n}\bigl(x_i^{(1)} - x_i^{(0)}\bigr)
  \left.\frac{\partial f_j}{\partial x_i}\right|_{x^{(0)}}
  = -f_j\bigl(x^{(0)}\bigr) .
\]
En évaluant ce système matriciel on obtient le vecteur $x_i^{(1)}$. On reprend
le processus et ainsi de suite. À chaque étape on résout le système
matriciel :
\[
  \sum_{i=1}^{n}\bigl(x_i^{(k+1)} - x_i^{(k)}\bigr)
  \left.\frac{\partial f_j}{\partial x_i}\right|_{x^{(k)}}
  = -f_j\bigl(x^{(k)}\bigr) .
\]
Dès que les quantités $x_i^{(k)}$ sont trouvées à la précision choisie, on a
alors l'ensemble des racines du système algébrique non linéaire.

\chapter{Traitement numérique des mesures expérimentales}

Ce chapitre est utile lorsque nous avons obtenu des mesures expérimentales et
voulons les traiter. Les techniques numériques utilisées pour le traitement de
ces données portent essentiellement sur l'interpolation et l'approximation des
fonctions.

\section{Généralités}

\subsection{Position du problème}

L'interpolation et l'approximation des fonctions sont des techniques
numériques très utilisées en sciences expérimentales. Le problème est le
suivant :

Nous avons une grandeur $y$ qui dépend d'une variable $x$. Nous connaissons
les valeurs $y_i$ de cette grandeur en un certain nombre de valeurs discrètes
ou points discrets $x_1, x_2, \dots, x_n$. Nous voulons approximer $y$ par une
fonction analytique $y = f(x)$ telle que
\[
  y_1 = f(x_1),\quad y_2 = f(x_2),\quad \dots,\quad y_n = f(x_n).
\]
L'expression de $f(x)$ pourra ainsi permettre de déterminer les valeurs de la
fonction en des points autres que ceux de la suite discrète. On parle alors de
l'art de lire entre les lignes d'une table ou d'\emph{interpolation}. Ces
points où on cherche les nouvelles valeurs de $y$ peuvent être inaccessibles
expérimentalement (soit parce qu'ils sont éloignés de la plage de mesure de
l'appareil utilisé, soit parce que la précision de l'appareil de mesure ne
permet pas d'atteindre ces points).

Un tel problème n'admet pas de solution univoquement déterminée. Notons aussi
que l'expression analytique écrite de $f(x)$ ne nous intéresse pas
particulièrement puisqu'elle peut être gérée par l'ordinateur pour lire entre
les lignes d'une table. Mais on peut avoir besoin d'une telle expression pour
les besoins d'analyse des courbes.

Le problème peut se poser dans des cas où la grandeur $y$ dépend d'une seule
variable $x$, de deux ou de trois variables.

\subsection{Type d'interpolation}

La fonction $f(x)$ peut avoir plusieurs formes. Lorsqu'elle est un polynôme,
on parle d'interpolation \textbf{polynomiale}. Lorsque $f(x)$ est une fonction
trigonométrique, on parle d'interpolation \textbf{trigonométrique}. On peut
aussi avoir une interpolation par les fonctions exponentielles, des polynômes
de Legendre, les fonctions de Bessel, etc. Chacune de ces formes est choisie
en fonction des prédictions dictées par la nature du problème et ce que l'on
attend.

Cependant en pratique, l'on choisit souvent la forme la plus simple pour
$f(x)$, notamment la forme polynomiale. Dans le cas où l'on sait que la
fonction est périodique, il est judicieux d'utiliser l'interpolation
trigonométrique. La justification de l'un ou de l'autre type d'interpolation
repose sur 2 théorèmes présentés par \textbf{Weierstrass} (en 1885).

\paragraph{\underline{Théorème 1}} Toute fonction continue dans un intervalle
$[a,b]$ peut être représentée dans cet intervalle et pour chaque degré de
précision par un polynôme
\[
  f(x) = \sum_{n=0}^{N} a_n x^n .
\]

\paragraph{\underline{Théorème 2}} Toute fonction continue de période $2\pi$
peut être représentée par une série trigonométrique limitée de la forme
\[
  f(x) = \sum_{n=0}^{N} \bigl(a_n\cos nx + b_n\sin nx\bigr).
\]
Les coefficients $a_n$ et $b_n$ sont déterminés par la connaissance des
valeurs discrètes de la fonction sur l'intervalle $[a,b]$.

\section{Interpolation polynomiale}

Un polynôme de degré $n$ est défini par ses $(n+1)$ coefficients. Ainsi un
polynôme d'interpolation de degré $n$ sera complètement déterminé par la
connaissance de $(n+1)$ valeurs discrètes $y_i$. Il existe plusieurs formules
pour déterminer l'interpolation polynomiale.

\subsection{Formule d'interpolation de Lagrange}

L'idée de base de la formule de Lagrange est de déterminer des polynômes
$f_k(x)$ de degré $n$ qui prennent les valeurs 1 au point $x_k$ et 0 en tout
point différent de $x_k$
\[
  \Longrightarrow\quad
  f_k(x_{k'}) =
  \begin{cases}
    1 & \text{si } k' = k\\
    0 & \text{si } k' \neq k
  \end{cases}
\]
Le polynôme d'interpolation de Lagrange aura alors la forme
\[
  f(x) = y_1f_1(x) + y_2f_2(x) + \dots + y_{n+1}f_{n+1}(x)
  \quad\Longrightarrow\quad
  f(x) = \sum_{k=1}^{n+1} y_kf_k(x).
\]

\paragraph{Détermination de $f_k(x)$}

Soit par exemple un ensemble de deux points $(x_1, y_1)$ et $(x_2, y_2)$. Les
polynômes $f_k$ sont de degré 1, c'est-à-dire
\[
  f_1(x) = a_1 + b_1x\,; \qquad f_2(x) = a_2 + b_2x
\]
\[
  \left\{
  \begin{aligned}
    f_1(x_1) &= 1\\
    f_1(x_2) &= 0
  \end{aligned}
  \right.
  \quad\text{et}\quad
  \left\{
  \begin{aligned}
    f_2(x_1) &= 0\\
    f_2(x_2) &= 1
  \end{aligned}
  \right.
  \quad\Longrightarrow\quad
  f_1(x) = \frac{x - x_2}{x_1 - x_2}, \qquad
  f_2(x) = \frac{x - x_1}{x_2 - x_1}.
\]
Donc $y = y_1f_1(x) + y_2f_2(x)$.

Si on avait un ensemble de trois points $(x_1, y_1)$, $(x_2, y_2)$,
$(x_3, y_3)$, alors les polynômes seront de degré 2 et définis par :
\[
  f_1(x) = \frac{(x - x_2)(x - x_3)}{(x_1 - x_2)(x_1 - x_3)}, \quad
  f_2(x) = \frac{(x - x_1)(x - x_3)}{(x_2 - x_1)(x_2 - x_3)}, \quad
  f_3(x) = \frac{(x - x_1)(x - x_2)}{(x_3 - x_1)(x_3 - x_2)} .
\]
Dans ce cas, on a $y = y_1f_1(x) + y_2f_2(x) + y_3f_3(x)$. On parle alors
d'\emph{interpolation quadratique}.

En généralisant pour un ensemble de $n+1$ points, les polynômes de degré $n$
tels que
\[
  f_k(x_{k'}) =
  \begin{cases}
    1 & \text{si } k' = k\\
    0 & \text{si } k' \neq k
  \end{cases}
\]
sont définis par
\[
  f_k(x) = \frac{\displaystyle\prod_{\substack{i=1\\ i\neq k}}^{n+1}(x - x_i)}
                {\displaystyle\prod_{\substack{i=1\\ i\neq k}}^{n+1}(x_k - x_i)}
  \qquad\text{et}\qquad
  f(x) = \sum_{k=1}^{n+1} f_k(x)\,y_k .
\]
Le polynôme $f(x)$ est unique pour un ensemble de points $x_i$.

\paragraph{Exemple} On donne le tableau suivant :
\begin{center}
  \begin{tabular}{|c|c|c|c|}
    \hline
    $x$ & $-1$ & $0$ & $1$\\ \hline
    $y$ & $4$ & $1$ & $2$\\ \hline
  \end{tabular}
\end{center}
$(x_1 = -1,\ x_2 = 0,\ x_3 = 1)$. D'après la formule de Lagrange, on a :
\[
  \begin{aligned}
    f_1(x) &= \frac{(x - x_2)(x - x_3)}{(x_1 - x_2)(x_1 - x_3)} = \frac{x(x-1)}{2}, &
    f_2(x) &= \frac{(x - x_1)(x - x_3)}{(x_2 - x_1)(x_2 - x_3)} = -(x^2 - 1),\\
    f_3(x) &= \frac{(x - x_1)(x - x_2)}{(x_3 - x_1)(x_3 - x_2)} = \frac{(x+1)x}{2}
  \end{aligned}
\]
d'où
\[
  \begin{aligned}
    f(x) &= y_1f_1(x) + y_2f_2(x) + y_3f_3(x)\\
    f(x) &= 2x(x-1) - (x^2 - 1) + x(x+1).
  \end{aligned}
\]

À partir de l'expression de la fonction de Lagrange, on peut évaluer les
valeurs $y_p$ pour toute valeur $x_p$ comprise entre deux points $x_k$
(interpolation) ou en dehors des points $x_k$ (extrapolation). Pour ce faire,
on utilise l'algorithme suivant :

\begin{algo}
  \State 1- Entrer les données $(x_k, y_k, x_p)$ avec $k = 1, 2, \dots, n+1$
  \State 2- $y_p \leftarrow 0$
  \State 3- \textbf{Pour} $k$ allant de $1$ jusqu'à $n+1$, \textbf{faire}
  \State \hspace{2.2em} $y_s \leftarrow 1$
  \State \hspace{2.2em} \textbf{Pour} $i$ allant de $1$ jusqu'à $n+1$, \textbf{faire}
  \State \hspace{4.4em} \textbf{si} $i \neq k$
  \State \hspace{5.5em} $y_s \leftarrow y_s\,\dfrac{(x_p - x_i)}{(x_k - x_i)}$
  \State \hspace{2.2em} \textbf{Fin pour}
  \State \hspace{2.2em} $y_p \leftarrow y_p + y_k\,y_s$
  \State \hspace{1em} \textbf{Fin pour}
  \State 4- \textbf{Écrire} $x_p$ et $y_p$
  \State \textbf{Fin}
\end{algo}

\begin{exercice}
Traduire cet algorithme dans un des langages de programmation de votre choix
(Pascal, Fortran, C).
\end{exercice}

\subsection{Formule d'interpolation de Newton}

\subsubsection{Par les différences divisées}

On appelle différence divisée première des variables $x_i$ et $x_j$
correspondant aux valeurs $y_i$ et $y_j$ la quantité $\delta(x_i, x_j)$
\[
  \delta(x_i, x_j) = \frac{y_i - y_j}{x_i - x_j} .
\]
La différence divisée seconde est
\[
  \delta(x_i, x_j, x_k) = \frac{\delta(x_i, x_j) - \delta(x_i, x_k)}{x_j - x_k} .
\]
La différence divisée d'ordre $m$ est
\[
  \delta(x_1, \dots, x_m) =
  \frac{\delta(x_1, \dots, x_{m-1}) - \delta(x_1, \dots, x_{m-2}, x_m)}{x_{m-1} - x_m} .
\]
Le polynôme d'interpolation de Newton se met sous la forme
\[
  \begin{aligned}
    y = y_1 &+ \delta(x_1, x_2)(x - x_1) + \delta(x_1, x_2, x_3)(x - x_1)(x - x_2) + \dots\\
            &\dots + \delta(x_1, x_2, \dots, x_{k+1})(x - x_1)(x - x_2)\cdots(x - x_k) .
  \end{aligned}
\]

\subsubsection{Par les différences progressives et régressives}

La quantité $\Delta y_k = y_{k+1} - y_k$ est appelée différence première
progressive. La différence seconde progressive est
$\Delta^2 y_k = \Delta y_{k+1} - \Delta y_k$ et la différence progressive
d'ordre $\alpha$ est définie par
\[
  \Delta^{\alpha} y_k = \Delta^{\alpha-1} y_{k+1} - \Delta^{\alpha-1} y_k .
\]
On définit également la différence régressive d'ordre $\alpha$ de la manière
suivante
\[
  \nabla^{\alpha} y_k = \nabla^{\alpha-1} y_k - \nabla^{\alpha-1} y_{k-1}
  \qquad (\nabla y_k = y_k - y_{k-1}).
\]
Lorsque l'on a $x_{k+1} = x_k + h = x_1 + kh$, la formule d'interpolation par
les différences divisées se réduit à :
\[
  f(x) = y_1 + \frac{\Delta y_1}{h}(x - x_1)
             + \frac{\Delta^2 y_1}{h^2\,2!}(x - x_1)(x - x_2) + \dots
             + \frac{\Delta^{k} y_1}{h^{k}\,k!}(x - x_1)(x - x_2)\cdots(x - x_k) .
\]
Cette dernière formule est souvent appelée \textbf{formule de
Newton-Gregory}.

Si par contre on a $x_{k+1} = x_k - h = x_1 - kh$, alors on peut utiliser la
formule de Newton-Gregory régressive en remplaçant dans la formule précédente
$\Delta$ par $\nabla$.

\subsection{Erreurs d'interpolation polynomiale}

Dans la méthode de Lagrange, en remplaçant la valeur exacte $y_{\exp}$ (que
l'on aurait obtenue expérimentalement) par $y_p$, on commet une erreur de
troncature qui peut être déterminée par le théorème suivant :

\paragraph{\underline{Théorème}} Soit $f(x)$ le polynôme de degré $n$ qui
satisfait la condition $f(x_k) = y_k$ $(k = 1, \dots, n+1)$ et définie dans
l'intervalle $[a,b]$ qui contient les points $x_k$. Si $f$ est de classe
$C^{n+1}$ sur $[a,b]$ ($n+1$ fois dérivable), alors pour toute valeur
$x_p \in [a,b]$, il existe une valeur $\xi \in [a,b]$ telle que :
\[
  y_{\exp} - y_p = \Bigl[f^{(n+1)}(\xi)\big/(n+1)!\Bigr]
                   (x_p - x_1)(x_p - x_2)\cdots(x_p - x_{n+1}) .
\]
Cette différence représente l'\textbf{erreur de troncature}.

En supposant que les dérivées de la fonction $f$ sont uniformément bornées sur
$[a,b]$, l'erreur de troncature \textbf{décroît lorsque $n$ croît}.

\subsection{Interpolation Spline}

\subsubsection{Problème et propriétés}

L'interpolation Spline a originellement une formulation en mécanique. Elle est
utilisée (en méthodes des éléments finis) pour résoudre numériquement les
équations aux dérivées partielles. Son importance relève des inconvénients de
l'interpolation de Lagrange.

Le problème associé à l'interpolation Spline est le suivant :

Soit $x_1, x_2, \dots, x_n$ points distincts compris dans un intervalle
$[a,b]$ et soit $f_1, f_2, \dots, f_n$ $(y_i = f_i)$ les valeurs de la fonction
en ces points, l'interpolation Spline est une interpolation cubique qui passe
par les points $(x_i, f_i)$. C'est une fonction $S$ qui admet les propriétés
suivantes :
\begin{enumerate}[label=(\alph*)]
  \item $S(x_i) = f_i$,
  \item $S(x)$, $S'(x)$, $S''(x)$ sont continues sur $[a,b]$,
  \item dans chaque intervalle $[x_i, x_{i+1}]$ pour $i = 1, \dots, n-1$,
        $S(x)$ est un polynôme cubique (qui peut avoir des coefficients
        différents dans chaque sous-intervalle),
  \item si $g(x)$ est une autre fonction qui satisfait les conditions (a), (b)
        et (c), alors on doit avoir
        \[
          \int_a^b \bigl(S''(x)\bigr)^2\,\dd x \;\le\; \int_a^b \bigl(g''(x)\bigr)^2\,\dd x .
        \]
\end{enumerate}
Dans la terminologie de la mécanique, ces propriétés peuvent s'énoncer de la
manière suivante :
\begin{enumerate}[label=(\alph*)]
  \item la courbe Spline doit passer par tous les n\oe uds ;
  \item la courbe Spline n'est pas brisée ou ne se courbe pas suivant un angle
        pointu ;
  \item entre deux n\oe uds, la courbe Spline est un polynôme de degré 3 ;
  \item la courbe Spline minimise l'énergie potentielle du système. Cette
        propriété implique que
        \[
          S''(a) = S''(b) = 0 .
        \]
\end{enumerate}

\subsubsection{Forme de l'interpolation Spline}

Puisque $S$ est un polynôme cubique pour $x_i \le x \le x_{i+1}$, $S''(x)$ est
alors linéaire et la formule d'interpolation de $S''(x)$ est :
\[
  S''(x) = S''(x_i) + \frac{x - x_i}{x_{i+1} - x_i}\bigl[S''(x_{i+1}) - S''(x_i)\bigr].
\]
Après intégration, on obtient
\[
  S'(x) = S'(x_i) + S''(x_i)(x - x_i)
        + \frac{S''(x_{i+1}) - S''(x_i)}{2(x_{i+1} - x_i)}(x - x_i)^2 .
\]
Une deuxième intégration donne :
\begin{equation}
  S(x) = f_i + S'(x_i)(x - x_i) + \frac{S''(x_i)}{2}(x - x_i)^2
       + \frac{S''(x_{i+1}) - S''(x_i)}{6(x_{i+1} - x_i)}(x - x_i)^3 .
  \tag{a}
\end{equation}
À partir de ces définitions, on peut établir les relations suivantes :
\begin{equation}
  S''_{i-1}h_{i-1} + 2(h_i + h_{i-1})S''_i + S''_{i+1}h_i
  = 6\left[\frac{f_{i+1} - f_i}{h_i} - \frac{f_i - f_{i-1}}{h_{i-1}}\right]
  \qquad (i = 2, \dots, n-1)
  \tag{b}
\end{equation}
où $h_i = x_{i+1} - x_i$, $S''_i = S''(x_i)$.

On a aussi pour $i = 2, \dots, n-1$
\[
  S'_i = \frac{f_{i+1} - f_i}{h_i} - S''_{i+1}\frac{h_i}{6} - S''_i\frac{h_i}{3} .
\]
Les valeurs $S''_1$ et $S''_n$ sont déterminées par la quatrième propriété de
l'interpolation Spline à savoir
\[
  S''(a) = S''(b) = 0 \quad\Longrightarrow\quad S''_1 = S''_n = 0 .
\]
L'équation (b) donne un système algébrique linéaire dont la résolution permet
de trouver tous les $S''_i$. On les substitue dans (a) pour avoir les
fonctions $S(x)$ dans les intervalles $(x_i\,;\,x_{i+1})$.

\subsection{Autres types d'interpolation polynomiale}

Dans les applications, d'autres renseignements peuvent être connus sur la
fonction $f$. Par exemple, si la parité est connue, alors le polynôme est écrit
sous la forme d'une somme de puissances paires de $x$ si la fonction est paire
ou de puissances impaires si la fonction est impaire. Les coefficients du
polynôme peuvent être déterminés comme dans le cas général.

De plus il peut aussi arriver qu'en plus des valeurs $y_k = f(x_k)$, nous
disposons des valeurs des dérivées $y'_k = f'(x_k)$. Dans ces cas on utilise
la formule d'interpolation de \textbf{Hermite}. Il existe également d'autres
types d'interpolation polynomiale telle que celle de \textbf{Stirling} définie
à partir des différences centrées sur base entière, et celle de
\textbf{Bessel} définie à partir des différences centrées sur base demi
entière.

La formule de Stirling s'utilise surtout lorsque $x_p$ est proche d'un n\oe ud
$x_k$ tandis que celle de Bessel est adaptée au cas où $x_p$ est proche de
\[
  x_{k+1/2} = \frac{x_k + x_{k+1}}{2} .
\]

\paragraph{Remarque} L'opérateur de différence centrée sur base entière est
défini par :
\[
  \delta f(x) = f\Bigl(x + \frac{h}{2}\Bigr) - f\Bigl(x - \frac{h}{2}\Bigr)
  \;\longrightarrow\;
  \delta f_k = f_{k+1/2} - f_{k-1/2} .
\]

\section{Interpolation trigonométrique, double et inverse}

\subsection{Interpolation trigonométrique}

Quand la fonction recherchée est périodique, on utilise l'interpolation
trigonométrique. Ce problème a été d'abord analysé par Gauss et ensuite par
Hermite.

La formule d'interpolation de Hermite est la suivante :
\[
  f_k(x) = \prod_{\substack{i=1\\ i\neq k}}^{n+1}\sin(x - x_i)
           \Bigg/ \prod_{\substack{i=1\\ i\neq k}}^{n+1}\sin(x_k - x_i)
  \qquad\text{et}\qquad
  f(x) = \sum_{k=1}^{n+1} f_k(x)\,y_k .
\]
La différence entre la formule de Gauss et celle de Hermite réside dans le
fait que chez Gauss les arguments de la fonction sinus sont divisés par deux.

\subsection{Interpolation double}

\subsubsection{Problème}

Au cours de certaines études, l'on peut être confronté à un problème
d'interpolation à deux arguments, notamment dans le cas des problèmes plans où
une grandeur physique $z$ dépend des coordonnées $x$ et $y$. Ce type de
problème peut se résoudre de deux manières :
\begin{enumerate}[label=(\roman*)]
  \item on peut d'abord interpoler par rapport à une variable, ensuite par
        rapport aux autres ;
  \item on peut aussi trouver une formule d'interpolation qui utilise les
        différences doubles.
\end{enumerate}

\subsubsection{Différences doubles ou à deux voies}

Soit la fonction $Z = z(x,y)$ dont la valeur en un point $(x_m, y_e)$ est
$z_{me}$. On note
\[
  \begin{aligned}
    \Delta^{10} z_{me} &= \Delta_x z_{me} = z_{m+1,e} - z_{m,e}\\
    \Delta^{01} z_{me} &= \Delta_y z_{me} = z_{m,e+1} - z_{me}
  \end{aligned}
\]
Par exemple
\[
  \begin{aligned}
    \Delta^{10} z_{11} &= z_{21} - z_{11}\\
    \Delta^{11} z_{11} &= \Delta^2_{xy} z_{11} = \Delta^{10}(z_{12} - z_{11})
                        = \Delta^{01} z_{21} - \Delta^{01} z_{11} .
  \end{aligned}
\]
De façon générale
\[
  \Delta^{\alpha\beta} z_{11} = \Delta^{\alpha 0} z_{1,\beta+1}
    - \beta\,\Delta^{\alpha 0} z_{1,\beta}
    + \frac{\beta(\beta-1)}{2}\,\Delta^{\alpha 0} z_{1,\beta-1}
    - \dots + (-1)^{\beta}\,\Delta^{\alpha 0} z_{11} .
\]

\subsubsection{Formule d'interpolation double}

Elle est donnée par
\[
  \begin{aligned}
    &z(x,y) = f(x,y) = z_{11}
      + \frac{x - x_1}{h}\Delta^{10} z_{11} + \frac{y - y_1}{l}\Delta^{01} z_{11}\\
    &\quad + \frac{1}{2!}\left[\frac{(x - x_1)(x - x_2)}{h^2}\Delta^{20} z_{11}
        + \frac{2(x - x_1)(y - y_1)}{hl}\Delta^{11} z_{11}
        + \frac{(y - y_1)(y - y_2)}{l^2}\Delta^{02} z_{11}\right]\\
    &\quad + \dots\\
    &\quad + \frac{1}{m!}\Biggl[\frac{(x - x_1)(x - x_2)\cdots(x - x_m)}{h^m}\Delta^{m0} z_{11}
        + \frac{m(x - x_1)\cdots(x - x_{m-1})(y - y_1)}{h^{m-1}l}\Delta^{(m-1)1} z_{11}\\
    &\qquad\quad + \frac{m(m-1)}{2}\,\frac{(x - x_1)\cdots(x - x_{m-2})(y - y_1)(y - y_2)}{h^{m-2}l^2}
                \Delta^{(m-2)2} z_{11} + \dots\\
    &\qquad\quad + \frac{(y - y_1)(y - y_2)\cdots(y - y_m)}{l^m}\Delta^{0m} z_{11}\Biggr]
  \end{aligned}
\]
où $h$ et $l$ sont les pas suivant $x$ et suivant $y$.

\subsection{Interpolation inverse}

Elle consiste à déterminer $x_p$ connaissant $y_p$. On peut dans ce cas
utiliser les formules d'interpolation décrites précédemment. Il suffira de
remplacer $y$ par $x$.

\section{Approximation par la méthode des moindres carrés}

\subsection{Position du problème}

La méthode d'approximation par les moindres carrés est utilisée dans de
nombreuses applications sous des noms différents : optimisation linéaire,
méthode de régression, lissage des courbes.

Les expériences en physique peuvent contenir des milliers de mesures.
L'utilisation des méthodes d'interpolation va alors conduire à des grandes
erreurs d'arrondi et des calculs assez longs. La méthode des moindres carrés
nous permet d'utiliser plusieurs fonctions pour obtenir une approximation
simple.

Par exemple, soit $f_1, f_2, \dots, f_m$ les valeurs correspondantes aux
points $x_1, x_2, \dots, x_m$ et soit
\[
  f(x) = a_0 + a_1x + a_2x^2 + \dots + a_nx^n
\]
un polynôme d'approximation de degré $n$ avec $n \ll m$.

Il est clair que puisque $n \ll m$, il n'est pas possible d'avoir $f$ telle
que $\forall k$, $f(x_k) = f_k$. Mais on peut choisir $f$ telle que
$f(x_k) \approx f_k$. La fonction $f$ sera bonne si la quantité (reste)
\[
  R = \bigl[f_1 - f(x_1)\bigr]^2 + \bigl[f_2 - f(x_2)\bigr]^2 + \dots
      + \bigl[f_m - f(x_m)\bigr]^2
\]
est la plus petite possible. Il s'agit donc de trouver la fonction $f$ qui
minimise le reste $R$.

\subsection{Résolution du problème de l'approximation par les moindres
carrés}

Pour résoudre ce problème, on doit établir un système d'équations linéaires de
la forme :
\[
  \frac{\partial R}{\partial a_k} = 0 \qquad\text{avec } k = 0, \dots, n .
\]
La solution de ce système permet alors de déterminer les coefficients $a_i$
de $f$.

Par exemple, considérons le cas polynomial de premier degré $(n = 1)$ où
\[
  f(x) = a_0 + a_1x .
\]
Le reste est
\[
  R = \sum_{k=1}^{m}\bigl[f_k - (a_0 + a_1x_k)\bigr]^2 .
\]
Pour que $R$ soit minimal, il faut que
$\dfrac{\partial R}{\partial a_0} = 0$ et
$\dfrac{\partial R}{\partial a_1} = 0$.

On trouve alors
\[
  a_0 = \frac{\sum x_k^2\sum f_k - \sum x_k\sum x_kf_k}{m\sum x_k^2 - \left(\sum x_k\right)^2}
  \qquad\text{et}\qquad
  a_1 = \frac{m\sum x_kf_k - \sum x_k\sum f_k}{m\sum x_k^2 - \left(\sum x_k\right)^2}
\]
(établir en exercice).

En général, on choisit $f(x)$ sous la forme
\[
  f(x) = a_0g_0(x) + a_1g_1(x) + \dots + a_ng_n(x)
\]
où les fonctions $g_k(x)$ peuvent être de natures différentes (polynôme,
fonction rationnelle, logarithme, exponentielle, etc.).

Le système d'équations à résoudre est alors défini de la manière suivante :
\[
  \left\{
  \begin{aligned}
    A_{00}a_0 + A_{01}a_1 + \dots + A_{0n}a_n &= b_0\\
    A_{10}a_0 + A_{11}a_1 + \dots + A_{1n}a_n &= b_1\\
    &\;\;\vdots\\
    A_{n0}a_0 + A_{n1}a_1 + \dots + A_{nn}a_n &= b_n
  \end{aligned}
  \right.
\]
où
\[
  b_k = \sum_{i=1}^{m} g_k(x_i)\,f_i, \qquad
  A_{kj} = \sum_{i=1}^{m} g_k(x_i)\,g_j(x_i)
  \qquad\text{avec}\quad
  \left\{
  \begin{aligned}
    j &= 0, 1, \dots, n\\
    k &= 0, 1, \dots, n
  \end{aligned}
  \right.
\]

\begin{exercice}[Exercice 1]
Écrire les algorithmes de calculs des coefficients $b_k$ et $A_{kj}$.
\end{exercice}

\begin{exercice}[Exercice 2]
Les résultats d'une expérience sont donnés dans le tableau suivant :
\begin{center}
  \begin{tabular}{|c|c|c|c|c|c|}
    \hline
    $x$ & 1 & 2 & 3 & 4 & 5\\ \hline
    $y$ & 1,6 & 2 & 2,3 & 2,4 & 2,5\\ \hline
  \end{tabular}
\end{center}
On se propose de déterminer l'expression $y = f(x)$. Pour cela, on choisit
\[
  \text{a)}\quad y = a_0 + \frac{a_1}{x}
  \qquad\text{et}\qquad
  \text{b)}\quad y = a_0 + \frac{a_1}{x} + a_2e^{x} .
\]
En utilisant la méthode des moindres carrés, calculer les coefficients $a_0$,
$a_1$ et $a_2$.
\end{exercice}

\begin{exercice}[Exercice 3]
Du tableau suivant utiliser la formule d'interpolation de Lagrange pour
déterminer $y_p$ pour $x_p = 102$ et $x_p$ pour $y_p = 13{,}5$.
\begin{center}
  \begin{tabular}{|c|c|c|c|c|c|}
    \hline
    $x$ & 93,0 & 96,2 & 100,0 & 104,2 & 108,7\\ \hline
    $y$ & 11,38 & 12,80 & 14,70 & 17,07 & 19,91\\ \hline
  \end{tabular}
\end{center}
\textit{Réponse :} $(x_p = 102$ et $y_p = 15{,}79)$, $(x_p = 97{,}65$ et
$y_p = 13{,}5)$.
\end{exercice}

\chapter{Intégration et dérivation numérique}

Dans certains problèmes scientifiques, on arrive souvent à des dérivées
premières ou d'ordre $n$ ou à des intégrales simples ou doubles qui ne peuvent
pas être calculées par des méthodes analytiques. Il faut alors faire appel à
des méthodes numériques qui permettent de trouver des solutions approchées en
utilisant des formes discrètes issues des définitions mathématiques de base.

Ce chapitre présente les méthodes de dérivation numérique (en partant des
formules classiques, des interpolations linéaires, quadratiques et cubiques)
et des méthodes d'intégration numériques (formules des trapèzes et de
Simpson).

\section{Dérivation numérique}

\subsection{Formules classiques}

Soit $f$ une fonction différentiable sur un intervalle $[a,b]$. Le nombre
dérivé de $f$ au point $x$ est défini par :
\[
  f'(x) = \lim_{h\to 0}\frac{f(x+h) - f(x)}{h} .
\]
On peut également utiliser la formule
\[
  f'(x) = \lim_{h\to 0}\frac{f(x + h/2) - f(x - h/2)}{h} .
\]
De la même manière, la dérivée seconde est donnée par :
\[
  f''(x) = \lim_{h\to 0}\frac{f(x+h) - 2f(x) + f(x-h)}{h^2} .
\]
Numériquement, on écrit
\[
  f'(x) = \frac{f(x+h) - f(x)}{h}
  \qquad\text{ou}\qquad
  f'(x) = \frac{f\bigl(x + \frac{h}{2}\bigr) - f\bigl(x - \frac{h}{2}\bigr)}{h}
\]
et
\[
  f''(x) = \frac{f(x+h) - 2f(x) + f(x-h)}{h^2}
\]
où $h$ est un nombre très petit : par exemple $h = 10^{-3}$, $10^{-1}$.

La formule de la dérivée première peut s'obtenir à partir de la formule de
l'interpolation linéaire. En effet, la formule de l'interpolation linéaire est
\[
  f(x) = f(x_1)\frac{x - x_2}{x_1 - x_2} + f(x_2)\frac{x - x_1}{x_2 - x_1}
       + \frac{1}{2}f''(\xi)(x - x_1)(x - x_2) .
\]
D'où
\[
  f'(x) = \frac{f(x_2) - f(x_1)}{x_2 - x_1}
        + \frac{1}{2}f''(\xi)\bigl[(x - x_1) + (x - x_2)\bigr]
        + \frac{1}{2}(x - x_1)(x - x_2)\frac{\dd}{\dd x}\bigl[f''\bigl(\xi(x)\bigr)\bigr] .
\]
L'erreur de troncature correspondant à la dérivée par interpolation linéaire
est donc
\[
  E = \frac{1}{2}f''(\xi)\bigl[(x - x_1) + (x - x_2)\bigr]
    + \frac{1}{2}(x - x_1)(x - x_2)\frac{\dd}{\dd x}\bigl[f''\bigl(\xi(x)\bigr)\bigr]
\]
avec $\xi \in [x_1\,;\,x_2]$.

Si nous posons $x = x_1$ et $x_2 = x_1 + h$, on obtient
\[
  f'(x_1) = \frac{f(x_1 + h) - f(x_1)}{h}
\]
avec une erreur de troncature $E = -\dfrac{1}{2}hf''(\xi)$.

C'est la formule classique donnée précédemment.

\subsection{Formules plus précises}

\subsubsection{Dérivée à partir de l'interpolation quadratique}

La formule d'interpolation quadratique est
\[
  f(x) = f(x_1)f_1(x) + f(x_2)f_2(x) + f(x_3)f_3(x)
\]
où
\[
  f_1(x) = \frac{(x - x_2)(x - x_3)}{(x_1 - x_2)(x_1 - x_3)}, \quad
  f_2(x) = \frac{(x - x_1)(x - x_3)}{(x_2 - x_1)(x_2 - x_3)}, \quad
  f_3(x) = \frac{(x - x_1)(x - x_2)}{(x_3 - x_1)(x_3 - x_2)} .
\]
En procédant comme dans le cas d'une interpolation linéaire, et en posant
$x = x_2$, on établit que :
\[
  f'(x) = \frac{f(x+h) - f(x-h)}{2h}
\]
avec une erreur de troncature
\[
  E = -\frac{1}{6}h^2f'''(\xi) \qquad\text{où } \xi \in [x - h\,;\,x + h] .
\]
Si l'on posait $x = x_1$, il viendrait le résultat suivant :

Si $f$ est de classe $C^4$ sur l'intervalle $[a,b]$, alors pour toutes valeurs
de $x$ et $h$ avec $a < x + 2h < b$, il existe $\xi \in [x\,;\,x + 2h]$ tel
que
\[
  f'(x) = \frac{1}{2h}\bigl[-3f(x) + 4f(x+h) - f(x+2h)\bigr] + \frac{1}{3}h^2f'''(\xi) .
\]

\subsubsection{Dérivée à partir de l'interpolation cubique}

On obtient le résultat suivant :

Si $f$ est de classe $C^6$ sur l'intervalle $[a,b]$ alors pour toute valeur de
$x$ et $h$ telles que $a < x - 2h < x + 2h < b$, il existe
$\xi \in [x - 2h\,;\,x + 2h]$ tel que
\[
  f'(x) = \frac{1}{12h}\bigl[f(x-2h) - 8f(x-h) + 8f(x+h) - f(x+2h)\bigr]
        + \frac{1}{30}h^4f^{(5)}(\xi) .
\]

\begin{exercice}[Exercice 1]
Établir la formule de dérivation numérique à partir de l'interpolation
quadratique.
\end{exercice}

\begin{exercice}[Exercice 2]
Au cours d'une expérience de dynamique, la position en fonction du temps est
donnée dans le tableau suivant :
\begin{center}
  \begin{tabular}{|c|c|c|c|c|c|c|}
    \hline
    $t$ & 1.0 & 1.10 & 1.20 & 1.30 & 1.40 & 1.50\\ \hline
    $x$ & 1.6487 & 1.7333 & 1.8221 & 1.9155 & 2.0138 & 2.1170\\ \hline
  \end{tabular}
\end{center}
Calculer les vitesses aux instants $t = 1.05$ ; $1.15$ ; $1.25$ ; $1.35$ et
$1.45$.
\end{exercice}

\section{Intégrations numériques -- Formules de Newton-Cotes}

\subsection*{L'intégrale définie}

\[
  I = \int_a^b f(x)\,\dd x
\]
représente un nombre qui, dans certains cas favorables peut être calculé
analytiquement. Dans le cas contraire, on doit faire appel au calcul
numérique. Pour ce faire, la fonction $f$ supposée continue sur $[a,b]$ est
remplacée par une formule d'interpolation. Généralement, $f$ est approximée
par un polynôme et les formules d'intégration qui s'en déduisent sont dites
\emph{formules de Newton-Cotes}. Ces formules supposent que les n\oe uds
d'interpolation sont équidistants.

\subsection{Méthode des trapèzes}

\subsubsection{Formule composite des trapèzes}

Dans la méthode des trapèzes, l'interpolation de la fonction $f$ est
linéaire. En supposant que l'intervalle $[a,b]$ est court, on a
\[
  f(x) = f(a)\frac{x - b}{a - b} + f(b)\frac{x - a}{b - a}
       + \frac{1}{2}f''(\xi)(x - a)(x - b) .
\]
Il vient alors
\[
  I \cong \frac{b - a}{2}\bigl[f(a) + f(b)\bigr] + E
\]
où $E$ est l'erreur de troncature définie par
\[
  E = \frac{1}{2}\int_a^b f''(\xi)(x - a)(x - b)\,\dd x .
\]
Pour évaluer $E$, nous utilisons le théorème de la moyenne pour les
intégrales.

\paragraph{Théorème de la moyenne} Soient $f$ et $g$ deux fonctions
continues sur $[a,b]$. Si $g(x)$ ne change pas de signe sur $[a,b]$, alors il
existe $\eta \in [a,b]$ tel que
\[
  \int_a^b f(x)g(x)\,\dd x = f(\eta)\int_a^b g(x)\,\dd x .
\]
Puisque le produit $(x - a)(x - b)$ reste négatif sur $[a,b]$, nous aurons
donc
\[
  E = -\frac{1}{12}f''(\xi)(b - a)^3 .
\]
Si $b - a = h$, alors
\[
  E = -\frac{1}{12}f''(\xi)h^3 .
\]
Lorsque la longueur de l'intervalle $[a,b]$ n'est pas assez petite, l'on
utilise la propriété d'additivité de l'intégrale pour faire un découpage en
morceaux réduits, à savoir :
\[
  \int_a^b f(x)\,\dd x = \sum_{k=1}^{n}\int_{x_{k-1}}^{x_k} f(x)\,\dd x
  \qquad\text{où } a = x_0 < x_1 < \dots < x_{n-1} < x_n = b .
\]
Prenons $x_k = x_0 + kh$, on obtient alors :
\[
  \int_a^b f(x)\,\dd x = \frac{h}{2}\bigl[f(x_0) + 2f(x_1) + 2f(x_2) + \dots
                         + 2f(x_{n-1}) + f(x_n)\bigr] + E
\]
où $E = \sum_{k=1}^{n} E_k$.

C'est la formule composite de la méthode des trapèzes.

Sachant que $E_k = -\dfrac{1}{12}h^3f''(\xi_k)$, on obtient
\[
  E = \sum_{k=1}^{n} E_k = -\frac{1}{12}h^3\sum_{k=1}^{n} f''(\xi_k)
    = -\frac{nh^3}{12}\,\frac{\sum_{k=1}^{n} f''(\xi_k)}{n}
    = -\frac{nh^3}{12}f''(\eta)
\]
où $\eta \in [a,b]$. Or $b - a = nh$. D'où
\[
  E = -\frac{h^2}{12}(b - a)f''(\eta) .
\]

\begin{exercice}[Exercice 3]
Calculer à la main $I = \displaystyle\int_1^2 \frac{\dd x}{x}$ pour $h = 0.2$,
puis pour $h = 0.1$. Comparer à la valeur exacte.
\end{exercice}

\subsubsection{Algorithme de la formule composite des trapèzes}

Cet algorithme permet de calculer une intégrale par la méthode composite des
trapèzes. Il utilise des tableaux pour $x_k$ et $f_k$.

\begin{algo}
  \State 1\textdegree- Définir la fonction $f$
  \State 2\textdegree- Donner $a$, $b$ et $n$
  \State 3\textdegree- $h \leftarrow \dfrac{b - a}{n}$, $S \leftarrow 0$
  \State 4\textdegree- \textbf{Pour} $k$ allant de $0$ jusqu'à $n$, \textbf{faire}
  \State \hspace{2.5em} $x_k \leftarrow a + kh$
  \State \hspace{2.5em} $f_k \leftarrow f(x_k)$
  \State \hspace{1.2em} \textbf{Fin pour}
  \State 5\textdegree- \textbf{Pour} $k$ allant de $1$ jusqu'à $n-1$, \textbf{faire}
  \State \hspace{2.5em} $S \leftarrow S + hf_k$
  \State \hspace{1.2em} \textbf{Fin pour}
  \State \hspace{1.2em} $S \leftarrow S + \dfrac{h}{2}(f_0 + f_n)$
  \State 6\textdegree- \textbf{Écrire} $S$
\end{algo}

\begin{exercice}[Exercice 4]
Écrire un autre algorithme qui n'utilise pas les tableaux.
\end{exercice}

\subsection{Méthode de Simpson}

Dans la méthode de Simpson, la fonction $f$ est approximée par un polynôme de
degré deux qui passe par les points $a$, $b$ et $c = \dfrac{a+b}{2}$. Cette
méthode permet d'obtenir :
\[
  \int_a^b f(x)\,\dd x = \frac{b - a}{6}\bigl[f(a) + 4f(c) + f(b)\bigr] + E
\]
avec
\[
  E = \frac{1}{6}\int_a^b f'''(\xi)(x - a)(x - b)(x - c)\,\dd x .
\]
En utilisant l'additivité des intégrales, on obtient :
\[
  \int_a^b f(x)\,\dd x = \frac{h}{3}\bigl[f(x_0) + 4f(x_1) + 2f(x_2) + 4f(x_3)
                         + 2f(x_4) + \dots + 4f(x_{n-1}) + f(x_n)\bigr] .
\]

\begin{exercice}[Exercice 5]
Écrire un algorithme pour la méthode de Simpson.
\end{exercice}

\begin{exercice}[Exercice 6]
Établir une formule d'intégration numérique qui utilise une interpolation de
degré 3 et qui interpole les points $a$, $b$, $c = \dfrac{a+b}{4}$ et
$d = \dfrac{3(a+b)}{4}$.
\end{exercice}

\begin{exercice}[Exercice 7]
Utiliser la méthode de Simpson pour calculer manuellement l'intégrale
$I = \displaystyle\int_0^1 \frac{4}{1+x^2}\,\dd x$ avec $h = 0.1$.
\end{exercice}

\section{Intégrales mixtes et doubles}

\subsection{Intégrale mixte}

C'est une intégrale de la forme
\[
  I = \int_a^b f(x)g(x)\,\dd x
\]
où l'expression mathématique de $g(x)$ est connue et où $f(x)$ est approchée
par une formule d'interpolation. Ce type d'intégrale intervient par exemple
dans le calcul des coefficients de Fourier du développement d'une fonction :
\[
  a_n = \frac{2}{T}\int_0^T f(x)\cos(nx)\,\dd x
  \qquad\text{et}\qquad
  b_n = \frac{2}{T}\int_0^T f(x)\sin(nx)\,\dd x .
\]
Dans ce cas $g(x) = \cos(nx)$ ou $\sin(nx)$. Ce type d'intégrale intervient
également dans la résolution des équations intégrales (par exemple
l'équation $f(x) = 5x + \int_0^1 (t + x)f(t)\,\dd t$).

Pour calculer ce type d'intégrale, on élabore des formules d'intégration de
type Newton-Cotes ($f(x)$ est interpolée par un polynôme).

\subsection{Interpolation double}

Il s'agit de calculer une intégrale de la forme
\[
  I = \int_{x=a}^{x=b}\int_{y=c}^{y=d} Z(x,y)\,\dd x\,\dd y .
\]
Pour ce faire, l'on peut utiliser la formule d'interpolation double de la
fonction $Z(x,y)$.

\chapter{Équations différentielles et équations intégro-différentielles}

Un ensemble très important des problèmes scientifiques est décrit par des
équations différentielles ordinaires et des équations aux dérivées partielles
linéaires et non linéaires, à conditions aux limites (ou aux frontières) ou à
conditions initiales. Dans certains cas, un tel problème peut se résoudre en
trouvant les solutions générales et particulières par les méthodes analytiques
et classiques (résolution directe, utilisation des facteurs intégrants,
utilisation de la transformée de Laplace, etc.).

Lorsqu'aucune méthode analytique ne permet de calculer la solution analytique,
il faut trouver les méthodes numériques pour approcher la valeur cherchée.

Ce chapitre est consacré à la présentation des méthodes numériques pour
résoudre les équations différentielles ordinaires à valeurs initiales et à
valeurs aux limites. On abordera aussi les méthodes de résolution numérique
des équations intégro-différentielles dans lesquelles il y a la présence des
dérivées de la variable inconnue et des intégrales de la variable inconnue.

\section{Généralités}

\subsection{Définition et classification}

Une équation différentielle est une relation qui lie une fonction $y$ à ses
dérivées. S'il y a une seule variable indépendante, on parle alors d'équation
différentielle ordinaire. Dans le cas de plusieurs variables indépendantes, on
parle d'équations aux dérivées partielles. Les équations différentielles
ordinaires de variable indépendante $x$ peuvent se mettre sous la forme
\begin{equation}
  F\bigl(x, y, y', \dots, y^{(n-1)}, y^{(n)}\bigr) = 0
  \tag{1}
\end{equation}
où $y^{(i)}$ est la dérivée d'ordre $i$ et $n$ est l'ordre de l'équation
différentielle et $a \le x \le b$.

Le degré d'une équation différentielle, si celle-ci peut être décrite comme un
polynôme, où les indéterminées sont des dérivées, est le degré de la dérivée de
l'ordre le plus élevé. Par exemple, l'équation
$(y'')^3 + (y')^5 + 7y = e^x$ est une équation différentielle de second ordre
et de degré 3. Si la fonction $F$ est un polynôme dont les degrés de $y$ et de
ses dérivées sont égaux à 0 ou 1 et si $F$ ne contient pas des produits de $y$
et de ses dérivées, ou des dérivées entre elles, alors l'équation
différentielle est dite \emph{linéaire}. Dans tous les autres cas, l'équation
est dite \emph{non linéaire}.

La résolution d'une équation différentielle d'ordre $n$ nécessite la
connaissance au préalable de $n$ constantes, notamment $n$ conditions imposées
à $y$ et à ses dérivées. Si l'on fixe la valeur de $y$ et de ses $(n-1)$
dérivées au point initial $x = a$, alors, on a un \emph{problème aux valeurs
initiales}. Si par contre les $n$ conditions sont imposées les unes en $x = a$
et les autres en $x = b$, on a un \emph{problème aux valeurs aux limites}.

\subsection{Origine des équations différentielles}

Les équations différentielles sont d'une importance fondamentale dans presque
tous les domaines de la science. Ceci est dû au fait que plusieurs lois et
phénomènes ont une formulation mathématique sous forme d'équations
différentielles. Par exemple, la seconde loi de Newton est une équation
différentielle du second ordre. De même la loi de la radioactivité est une
équation du premier ordre.

\subsection{Position du problème numérique}

Considérons l'équation différentielle ordinaire du premier ordre
\begin{equation}
  y' = f(x, y) .
  \tag{2}
\end{equation}
Le problème est de trouver la valeur $\bar{y}$ de $y$ correspondant à
$x = x_0 + h$ (ou en $x_k = x_0 + kh$) connaissant $y = y_0$ pour $x = x_0$.
$h$ est un nombre réel et $k$ un nombre entier.

En intégrant (2), on obtient
\begin{equation}
  y = y_0 + \int_{x_0}^{x} f(x, y)\,\dd x .
  \tag{3}
\end{equation}
On a alors
\begin{equation}
  \bar{y} = y_0 + \int_{x_0}^{x_0+h} f(x, y)\,\dd x = y_0 + L
  \qquad\text{où}\qquad
  L = \int_{x_0}^{x_0+h} f(x, y)\,\dd x .
  \tag{4}
\end{equation}
On peut donc utiliser des méthodes numériques de calcul des intégrales pour
calculer l'intégrale. On peut également utiliser le développement de
$y(x+h)$ en série de puissances de $h$ (puis insérer les différentes dérivées
à partir de l'équation (2)).

\section{Méthodes pour les problèmes à valeurs initiales}

\subsection{Méthode d'Euler}

Encore appelée méthode de la dérivée première, elle s'exprime par la formule
\[
  y(x+h) = y(x) + hf(x, y) .
\]
La formule itérative est donc
\[
  y_{k+1} = y_k + hf(x_k, y_k) .
\]
C'est la méthode la plus simple. Mais son inconvénient réside dans le fait
qu'elle est souvent très instable et exige des pas de calcul $h$ très petits.

\paragraph{Algorithme de la méthode d'Euler}

\begin{algo}
  \State Définir $f(x,y)$
  \State $x \leftarrow x_0$, $y \leftarrow y_0$
  \State Donner la valeur de $h$
  \For{$k$ allant de $1$ à $n$}
    \State $y \leftarrow y + hf(x,y)$
    \State $x \leftarrow x + h$
    \Print $x$ et $y$
  \EndFor
\end{algo}

\subsection{Les méthodes de Runge et de Kutta}

Ce sont les méthodes élaborées par C.~Runge en 1894 et améliorées par
W.~Kutta vers 1901. Elles utilisent les méthodes d'intégration numérique des
trapèzes et de Simpson. Ces méthodes sont les plus utilisées car elles sont
très stables.

\subsubsection{Méthode de Runge-Kutta du second ordre}

Elle est basée sur la formule intégrale des trapèzes (voir chapitre sur
l'intégration numérique) et son schéma itératif se met sous la forme :
\[
  y_{k+1} = y_k + \frac{h}{2}\Bigl\{f(x_k, y_k) + f\bigl[x_k + h,\; y_k + hf(x_k, y_k)\bigr]\Bigr\} .
\]
Elle présente une erreur d'ordre $O(h^3)$.

\begin{exercice}
Écrire l'algorithme de la méthode de Runge-Kutta d'ordre 2.
\end{exercice}

\subsubsection{Méthode de Runge du quatrième ordre}

Supposons les valeurs $y_0$, $y_1$, $y_2$ de $y$ connues aux points $x_0$,
$x_1 = x_0 + \dfrac{h}{2}$ et $x_2 = x_0 + h$.

En utilisant la formule d'intégration numérique de Simpson (voir chapitre
précédent), on obtient
\[
  L = \frac{h}{6}\Bigl[f(x_0, y_0) + 4f\Bigl(x_0 + \frac{h}{2}, y_1\Bigr) + f(x_0 + h, y_2)\Bigr].
\]
Les valeurs $y_1$ et $y_2$ étant inconnues, on utilise les approximations
\[
  y_1 \approx y_0 + \frac{h}{2}f(x_0, y_0) = y_0 + \frac{h}{2}f_0, \qquad
  y_2 \approx y_0 + hf(x_0 + h, y_0 + hf_0) .
\]
Ces approximations sont choisies de telle manière qu'en développant $L$ en
séries de puissance de $h$, ces 3 premiers termes coïncident avec le
développement en série de Taylor de $y(x+h)$. Il vient alors
\[
  L \approx \frac{h}{6}\Bigl\{f(x_0, y_0)
     + 4f\Bigl[x_0 + \frac{h}{2},\; y_0 + \frac{h}{2}f(x_0, y_0)\Bigr]
     + f\Bigl[x_0 + h,\; y_0 + hf\bigl(x_0 + h, y_0 + hf(x_0, y_0)\bigr)\Bigr]\Bigr\}.
\]
Soit plus généralement,
\[
  y_{k+1} = y_k + L
\]
avec
\[
  L = \frac{h}{6}\Bigl\{f(x_k, y_k)
     + 4f\Bigl[x_k + \frac{h}{2},\; y_k + \frac{h}{2}f(x_k, y_k)\Bigr]
     + f\Bigl[x_k + h,\; y_k + hf\bigl(x_k + h, y_k + hf(x_k, y_k)\bigr)\Bigr]\Bigr\}.
\]
Elle présente une erreur d'ordre $O(h^5)$.

Pratiquement on fait le calcul de la manière suivante. On calcule les quantités
$L_i$ définies par :
\[
  \begin{aligned}
    L_1 &= hf(x_k, y_k)\\
    L_2 &= hf(x_k + h, y_k + L_1)\\
    L_3 &= hf(x_k + h, y_k + L_2)\\
    L_4 &= hf\Bigl(x_k + \frac{h}{2}, y_k + \frac{L_1}{2}\Bigr)
  \end{aligned}
\]
Par la suite, on écrit
\[
  y_{k+1} = y_k + \frac{1}{6}\bigl(L_1 + 4L_4 + L_3\bigr).
\]

\begin{exercice}
Établir l'algorithme de la méthode de Runge d'ordre 4.
\end{exercice}

\subsubsection{Méthode de Runge-Kutta ou de Kutta-Simpson du quatrième ordre}

C'est une amélioration de la méthode de Runge d'ordre 4. Elle présente
également une erreur d'ordre $O(h^5)$ et son schéma itératif s'écrit :
\[
  y_{k+1} = y_k + \frac{1}{6}\bigl(L_1 + 2L_2 + 2L_3 + L_4\bigr)
\]
avec
\[
  \begin{aligned}
    L_1 &= hf(x_k, y_k)\\
    L_2 &= hf\Bigl(x_k + \frac{h}{2}, y_k + \frac{L_1}{2}\Bigr)\\
    L_3 &= hf\Bigl(x_k + \frac{h}{2}, y_k + \frac{L_2}{2}\Bigr)\\
    L_4 &= hf(x_k + h, y_k + L_3)
  \end{aligned}
\]
C'est la méthode la plus utilisée. Il existe d'autres variantes des méthodes de
Runge-Kutta telles que la méthode de Kutta-Merson ou de Kutta-Simpson du
sixième ordre.

\subsection{Méthodes à pas multiples ou d'Adams-Bashforth}

\subsubsection{Formule du second ordre}

Soit l'équation différentielle
\[
  y' = f(x, y) .
\]
On veut connaître $y(x+h)$ connaissant $y(x)$ et $y(x-h)$. Pour cela, on fait
le développement suivant :
\[
  y(x+h) = y(x) + hy' + \frac{h^2}{2}y'' + O(h^3)
  \quad\Longrightarrow\quad
  y_{k+1} = y_k + hy'_k + \frac{h^2}{2}y''_k .
\]
On a aussi
\[
  y'(x-h) = y'(x) - hy''(x) .
\]
D'où, il vient alors
\[
  y''_k = \frac{y'_k - y'_{k-1}}{h} .
\]
Finalement, on obtient
\[
  y_{k+1} = y_k + \frac{h}{2}\bigl[3y'_k - y'_{k-1}\bigr]
  \quad\Longrightarrow\quad
  y_{k+1} = y_k + \frac{h}{2}\bigl[3f(x_k, y_k) - f(x_{k-1}, y_{k-1})\bigr].
\]
C'est la formule d'Adams-Bashforth du second ordre présentant une erreur
d'ordre $O(h^3)$.

Pratiquement au début des itérations, il faut connaître les valeurs de $y$ en
$x_0$ et en $x_1$ afin de calculer la première valeur $y_2$ de l'itération.
Pour avoir $y_1$, on peut utiliser l'une des formules approchées classiques.
Par exemple en utilisant la formule d'intégration des trapèzes, on obtient :
\[
  y_1 = y_0 + \frac{h}{2}\Bigl[f(x_0, y_0) + f\bigl[x_0 + h, y_0 + hf(x_0, y_0)\bigr]\Bigr].
\]
La formule itérative peut alors débuter par $k = 1$, et on calcule les $y_2$,
$y_3$, $y_4$\dots{} Par exemple,
\[
  y_2 = y_1 + \frac{h}{2}\bigl[3f(x_1, y_1) - f(x_0, y_0)\bigr].
\]

\begin{exercice}
Écrire l'algorithme de la méthode d'Adams-Bashforth du second ordre.
\end{exercice}

\subsubsection{Formules d'ordre supérieur}

Pour obtenir des formules d'ordre supérieur présentant par exemple des erreurs
d'ordre $O(h^4)$ ou $O(h^5)$, il suffit de développer à l'ordre supérieur
$y(x+h)$ et de remplacer les différentes dérivées.

\subsection{Méthode de prédicteur-correcteur}

Les méthodes de Runge, Runge-Kutta et de Kutta-Simpson sont fondamentalement
implicites car l'inconnue $y_{k+1}$ est exprimée à l'aide d'une fonction
contenant $y_{k+1}$. Par exemple, dans le cas de Runge-Kutta d'ordre 2, on a
normalement
\[
  y_{k+1} = y_k + \frac{h}{2}\bigl[f(x_k, y_k) + f(x_{k+1}, y_{k+1})\bigr].
\]
Afin d'avoir une expression explicite, on part des approximations pour
exprimer $y_{k+1}$ du second membre de l'égalité en fonction de $y_k$. On
pourrait aussi partir des formules d'Adams-Bashforth pour évaluer (prédire) les
valeurs de $y_{k+1}$, puis insérer ces valeurs prédites dans le second membre
de la formule implicite afin d'avoir une valeur corrigée. Dans un tel
processus, la formule explicite est appelée \emph{prédicteur} tandis que la
formule implicite utilisée est appelée \emph{correcteur}. Par exemple une
méthode de prédicteur-correcteur peut utiliser les formules suivantes :
\begin{align*}
  y_{k+1} &= y_k + \frac{h}{2}\bigl[3f(x_k, y_k) - f(x_{k-1}, y_{k-1})\bigr]\\
          &\equiv \text{prédicteur (formule d'Adams-Bashforth du second ordre),}\\[0.6em]
  y_{k+1} &= y_k + \frac{h}{2}\bigl[f(x_k, y_k) + f(x_{k+1}, y_{k+1})\bigr]\\
          &\equiv \text{correcteur (Runge-Kutta d'ordre 2).}
\end{align*}
On peut également combiner les autres formules implicites de RK aux autres
formules explicites d'Adams-Bashforth pour avoir d'autres méthodes de
prédicteur-correcteur plus élaborées.

\begin{exercice}
Établir l'algorithme de la méthode de prédicteur-correcteur, utilisant comme
prédicteur la formule d'Adams-Bashforth du second ordre et comme correcteur la
formule de $RK_2$ et traduire cet algorithme en Pascal ou en Fortran ou en C.
\end{exercice}

\section{Méthodes pour les problèmes aux valeurs limites}

Il s'agit des problèmes de la forme
\[
  y^{(n)} = f\bigl(x, y, y', \dots, y^{(n-1)}\bigr), \qquad a < x < b
\]
avec $y(a) = y_1$ et $y(b) = y_2$ et d'autres conditions sur les dérivées. On
peut résoudre ce type de problème par plusieurs méthodes dont la méthode de
tir, la méthode des fonctions complémentaires ou de superposition et la
méthode des différences finies.

\subsection{Méthode de tir}

Elle est utilisée dans le cas des problèmes aux limites non linéaires. Il
s'agit de transformer un problème aux valeurs aux limites en un problème aux
valeurs initiales par un choix adéquat des valeurs initiales. Ce choix se fait
à l'une des bornes du domaine d'intégration et doit être tel que la solution
passe le plus près possible de la valeur de la grandeur à l'autre borne du
domaine d'intégration.

Par exemple, considérons une équation du second ordre
\[
  y'' = f(x, y, y') \quad\text{avec}\quad y(a) = \eta_1 \quad\text{et}\quad y(b) = \eta_2 .
\]
La méthode de tir consiste à choisir $\gamma$ telle que $y'(a) = \gamma$ et de
résoudre le problème aux valeurs initiales
\[
  y'' = f(x, y, y') \quad\text{avec}\quad y(a) = \eta_1 \quad\text{et}\quad y'(a) = \gamma .
\]
Le choix de $\gamma$ est bon lorsqu'en $x = b$, on a $y(b) \approx \eta_2$.
Sinon il faut choisir une autre valeur pour $\gamma$. Une procédure à utiliser
pour le choix de $\gamma$ est la suivante. On sait que la valeur $y(b)$ dépend
de $\gamma$ ; c'est-à-dire que $y(b) = y(b, \gamma)$. Pour un bon choix de
$\gamma$, on doit avoir
\[
  y(b, \gamma) - \eta_2 \approx 0 \;\Longleftrightarrow\; f(\gamma) = 0
  \quad\text{où}\quad f(\gamma) = y(b, \gamma) - \eta_2 .
\]
Or d'après la méthode de Newton (voir chapitre II), le schéma itératif donnant
la solution de cette équation algébrique est
\[
  \gamma_{n+1} = \gamma_n - \frac{f(\gamma_n)}{f'(\gamma_n)}
  = \gamma_n - \frac{y(b, \gamma_n) - \eta_2}{y'(b, \gamma_n)}
  = \gamma_n - \frac{(\gamma_n - \gamma_{n-1})\bigl(y(b, \gamma_n) - \eta_2\bigr)}
                    {y(b, \gamma_n) - y(b, \gamma_{n-1})}
\]

Ainsi en pratique, on commence par donner une valeur initiale $\gamma_0$ à
$\gamma$. Ensuite, on choisit la valeur suivante $\gamma_1$, les autres valeurs
de $\gamma$ $(\gamma_2, \gamma_3, \dots)$ sont ensuite évaluées à partir de la
formule ci-dessus.

\begin{exercice}
Établir l'algorithme de la méthode de tir.
\end{exercice}

\subsection{Méthode des fonctions complémentaires ou de superposition}

C'est une méthode qui transforme une équation différentielle linéaire à
valeurs aux limites en deux problèmes différentiels à valeurs initiales. Elle a
plusieurs variantes que nous exposerons sur deux équations différentielles du
second ordre.

\subsubsection{Première variante}

Considérons l'équation
\[
  \left\{
  \begin{aligned}
    &y'' = p(x)y' + q(x)y + r(x)\\
    &y(a) = \eta_1 \quad\text{et}\quad y(b) = \eta_2
  \end{aligned}
  \right.
\]
Supposons que :
\[
  y(x) = y_1(x) + \mu y_2(x) \qquad (\mu = \text{cte} \neq 0).
\]
Il vient
\[
  \begin{aligned}
    y''_1 &= p(x)y'_1 + q(x)y_1 + r(x)\\
    y''_2 &= p(x)y'_2 + q(x)y_2
  \end{aligned}
\]
En $x = a$, on a $y_1(a) + \mu y_2(a) = \eta_1$.

Imposons que $y_1(a) = \eta_1 \Longrightarrow y_2(a) = 0$.

En $x = b$, on a $y_1(b) + \mu y_2(b) = \eta_2$. D'où
\[
  \Longrightarrow\quad \mu = \frac{\eta_2 - y_1(b)}{y_2(b)} .
\]
Fixons
\[
  y'_1(a) = 0 \qquad\text{et}\qquad y'_2(a) = 1 .
\]
On obtient alors les deux problèmes différentiels à valeurs initiales
suivants :
\[
  \left\{
  \begin{aligned}
    &\left\{
    \begin{aligned}
      &y''_1 = p(x)y'_1 + q(x)y_1 + r(x)\\
      &y_1(a) = \eta_1, \quad y'_1(a) = 0
    \end{aligned}
    \right.\\[0.4em]
    &\left\{
    \begin{aligned}
      &y''_2 = p(x)y'_2 + q(x)y_2\\
      &y_2(a) = 0, \quad y'_2(a) = 1
    \end{aligned}
    \right.
  \end{aligned}
  \right.
  \qquad x \in\, ]a, b]
\]
On résout alors numériquement ces deux problèmes à valeurs initiales. On
obtient les valeurs de $y_1(b)$ et $y_2(b)$ et par conséquent la valeur de
$\mu$ et les différentes valeurs de $y$ aux n\oe uds de discrétisation
(puisque les valeurs de $y_1$ et $y_2$ sont connues en tout point de
l'intervalle $]a,b]$).

\subsubsection{Deuxième variante}

Considérons le problème suivant
\begin{align}
  y'' &= p(x)y + q(x) \tag{a}\\
  y'(a) &= \alpha_{00}y(a) + \alpha_{10} \tag{b}\\
  y'(b) &= \beta_{00}y(b) + \beta_{10} \tag{c}
\end{align}
Considérons également l'équation
\begin{equation}
  y' = \alpha_0(x)y + \alpha_1(x) \tag{d}
\end{equation}
où nous avons choisi $\alpha_0(x)$ et $\alpha_1(x)$ tels que $y$ satisfait
toujours l'équation (a). En différenciant l'équation (d) par rapport à $x$ et
en identifiant à l'équation (a), on obtient
\begin{align}
  \alpha'_0 + \alpha_0^2 &= p(x) \tag{e}\\
  \alpha'_1 + \alpha_0\alpha_1 &= q(x) \tag{f}
\end{align}
avec $\alpha_0(a) = \alpha_{00}$, $\alpha_1(a) = \alpha_{10}$.

Les équations (e) et (f) sont des équations différentielles du
1\textsuperscript{er} ordre que l'on peut résoudre analytiquement ou
numériquement pour trouver les valeurs de $\alpha_0(x)$ et de $\alpha_1(x)$.

De l'équation (d), on a
$y'(b) = \alpha_0(b)y(b) + \alpha_1(b) = \beta_{00}y(b) + \beta_{10}$. On
obtient alors
\begin{align}
  y(b) &= \frac{\beta_{10} - \alpha_1(b)}{\alpha_0(b) - \beta_{00}} \tag{g}\\
  y'(b) &= \frac{\beta_{00}\alpha_1(b) - \beta_{10}\alpha_0(b)}{\beta_{00} - \alpha_0(b)} \tag{h}
\end{align}
L'équation (a) est ainsi transformée en un problème à valeurs initiales. On
peut en effet résoudre le problème $y'' = p(x)y + q(x)$ avec les conditions
initiales (g) et (h) ou résoudre le problème $y' = \alpha_0(x)y + \alpha_1(x)$
avec la condition initiale (g).

\subsection{Méthode des différences finies}

Elle convertit le problème de résolution d'une équation différentielle à la
résolution des équations algébriques.

Soit le problème :
\[
  y'' = p(x)y' + q(x)y + r(x) \qquad\text{avec}\quad y(a) = \eta_1 \quad\text{et}\quad y(b) = \eta_2 .
\]
On a :
\[
  y'(x) = \frac{y(x+h) - y(x)}{h}
  \qquad\text{et}\qquad
  y''(x) = \frac{y(x+h) - 2y(x) + y(x-h)}{h^2} .
\]
Soit $x = x_k$, on posera
\[
  y(x_k) = y_k, \quad p(x_k) = p_k, \quad q(x_k) = q_k, \quad r(x_k) = r_k .
\]
Par substitution des expressions des dérivées dans l'équation différentielle,
on aboutit à l'équation discrète
\[
  (1 - hp_k)\,y_{k+1} - \bigl(2 - hp_k + h^2q_k\bigr)\,y_k + y_{k-1} = h^2r_k
\]
pour $k = 1$ à $N-1$ et $h = \dfrac{b - a}{N}$.

On aboutit ainsi à un système d'équations algébriques linéaires de la forme
\[
  \left\{
  \begin{aligned}
    &y_0 = \eta_1\\
    &(1 - hp_1)\,y_2 - \bigl(2 - hp_1 + h^2q_1\bigr)\,y_1 + y_0 = h^2r_1\\
    &(1 - hp_2)\,y_3 - \bigl(2 - hp_2 + h^2q_2\bigr)\,y_2 + y_1 = h^2r_2\\
    &\qquad\vdots\\
    &(1 - hp_{N-1})\,y_N - \bigl(2 - hp_{N-1} + h^2q_{N-1}\bigr)\,y_{N-1} + y_{N-2} = h^2r_{N-1}\\
    &y_N = \eta_2
  \end{aligned}
  \right.
\]
Ce système algébrique peut être résolu par les méthodes itératives de Jacobi
ou de Gauss-Seidel.

\section{Les équations intégrales et intégro-différentielles}

\subsection{Définition}

On appelle équation intégrale une équation dont la fonction inconnue se trouve
dans le symbole d'intégration.

C'est l'exemple de l'équation suivante :
\[
  y(x) = e^x\sin x + 2\int_0^x \cos(x - t)\,y(t)\,\dd t .
\]
On a aussi les équations intégro-différentielles comme par exemple
\[
  y'(x) = 3y(x) + 2x + \int_0^x \sin w(x - t)\,y(t)\,\dd t .
\]
Une équation intégrale est dite linéaire lorsque la fonction inconnue $y$
intervient avec un degré 1. Sous la forme générale, les équations intégrales
linéaires se présentent de la manière suivante
\[
  y(x) = f(x) + \int_0^x k(x,t)\,y(t)\,\dd t .
\]
La fonction $k(x,t)$ est appelée \emph{noyau} de l'équation.

Dans les cas simples, les équations intégrales peuvent se résoudre
analytiquement en utilisant par exemple la méthode de la fonction de Gauss, de
la transformée de Laplace, de la transformation de Mellin, ou par les
approximations successives. Dans le cas contraire, il faut utiliser les
méthodes numériques.

\subsection{Résolution numérique des équations intégrales linéaires}

Soit l'équation linéaire
\begin{equation}
  y(x) = f(x) + \int_a^b k(x,t)\,y(t)\,\dd t .
  \tag{1}
\end{equation}
L'intégrale définie dans (1) peut être remplacée par les formules des trapèzes
et de Simpson. On obtient alors
\begin{equation}
  y(x) = f(x) + (b - a)\bigl[c_1k(x,t_1)y(t_1) + c_2k(x,t_2)y(t_2) + \dots
         + c_nk(x,t_n)y(t_n)\bigr]
  \tag{2}
\end{equation}
où $t_1, t_2, \dots, t_n$ sont les n\oe uds de subdivision de l'intervalle
$[a,b]$ et les coefficients $c_i$ dépendent de la formule d'intégration
utilisée.

Puisque l'équation (2) est vérifiée pour tout $x \in [a,b]$, elle est également
vérifiée pour $x = t_1$, $x = t_2$, \dots, $x = t_n$. Par conséquent, à partir
de l'équation (2), on obtient un système d'équations de la forme
\begin{equation}
  y(t_i) = f(t_i) + (b - a)\bigl[c_1k(t_i,t_1)y(t_1) + c_2k(t_i,t_2)y(t_2) + \dots
           + c_nk(t_i,t_n)y(t_n)\bigr]
  \tag{3}
\end{equation}
pour $i$ allant de $1$ à $n$ ($t_1 = a$, \dots, $t_i = a + (i-1)h$, \dots,
$t_n = b$).

Posons
\[
  y(t_i) = y_i \qquad\text{et}\qquad f(t_i) = f_i .
\]
Le système (3) devient alors
\begin{equation}
  \left\{
  \begin{aligned}
    y_1 &= f_1 + (b - a)\bigl[c_1k(t_1,t_1)y_1 + c_2k(t_1,t_2)y_2 + \dots + c_nk(t_1,t_n)y_n\bigr]\\
    y_2 &= f_2 + (b - a)\bigl[c_1k(t_2,t_1)y_1 + c_2k(t_2,t_2)y_2 + \dots + c_nk(t_2,t_n)y_n\bigr]\\
    &\;\;\vdots\\
    y_n &= f_n + (b - a)\bigl[c_1k(t_n,t_1)y_1 + c_2k(t_n,t_2)y_2 + \dots + c_nk(t_n,t_n)y_n\bigr]
  \end{aligned}
  \right.
  \tag{4}
\end{equation}
C'est un système de $n$ équations à $n$ inconnues que sont les $y_i$. La
résolution de ce système permet de déterminer les $y_i$ et en les substituant
dans l'équation (2), on obtient l'expression analytique de $y(x)$, solution de
l'équation (1).

\paragraph{Remarque} On peut transformer une équation différentielle
ordinaire en une équation intégrale. En résolvant l'équation intégrale, on
trouve alors une solution analytique de l'équation différentielle.

\paragraph{Exemple} Soit l'équation
\[
  y(x) = \frac{5x}{6} - \frac{1}{9} + \frac{1}{3}\int_0^1 (t + x)\,y(t)\,\dd t .
\]
Ici, on a $f(x) = 5x/6 - 1/9$ et $k(x,t) = x + t$.

Pour résoudre numériquement cette équation, utilisons la méthode de Simpson et
considérons pour cela trois points : $t_1 = 0$, $t_2 = 0{,}5$ et $t_3 = 1$. Le
système algébrique à résoudre est alors défini par :
\[
  y(x) = \frac{5x}{6} - \frac{1}{9}
       + \frac{1}{3}\times\frac{1}{3}\times\frac{1}{2}
         \Bigl[(t_1 + x)y(t_1) + 4(t_2 + x)y(t_2) + (t_3 + x)y(t_3)\Bigr]
\]
car $h = 1/2$.

En remplaçant $x$ par $t_i$, on obtient
\[
  \left\{
  \begin{aligned}
    y_1 &= -\frac{1}{9} + \frac{1}{18}\bigl(2y_2 + y_3\bigr) && \text{en } t = 0\\
    y_2 &= \frac{5}{12} - \frac{1}{9} + \frac{1}{18}\Bigl(\frac{y_1}{2} + 4y_2 + \frac{3}{2}y_3\Bigr) && \text{en } t = 0{,}5\\
    y_3 &= \frac{5}{6} - \frac{1}{9} + \frac{1}{18}\bigl(y_1 + 6y_2 + 2y_3\bigr) && \text{en } t = 1
  \end{aligned}
  \right.
\]
On trouve $y_1 = 0$, $y_2 = \frac{1}{2}$ et $y_3 = 1$. En remplaçant
$y_1 = y(t_1)$, $y_2 = y(t_2)$ et $y_3 = y(t_3)$ dans l'expression de $y(x)$,
on obtient $y(x) = x$.

\begin{exercice}
Discrétiser l'équation intégro-différentielle suivante :
\[
  y'(x) = 2y(x) + f(x) + \int_0^x k(x,t)\,y(t)\,\dd t .
\]
\end{exercice}

\chapter{Équations aux dérivées partielles}

Les équations aux dérivées partielles décrivent de nombreux problèmes de
physique et de mécanique. Leur résolution numérique fait appel à des formules
de discrétisation qui convertissent l'équation aux dérivées partielles en un
système d'équations discrètes que l'on peut résoudre par des schémas
itératifs.

Ce chapitre présente les méthodes de base de discrétisation des équations aux
dérivées partielles. En particulier, on s'intéresse ici aux équations de
Poisson et de Laplace, à l'équation de la chaleur et à l'équation d'ondes.
Pour que les schémas discrets conduisent à des solutions viables, ils doivent
être stables, c'est-à-dire que les solutions numériques doivent être très peu
sensibles au changement des pas de discrétisation. C'est pourquoi, ce chapitre
s'intéresse aussi à l'établissement des conditions de stabilité des schémas
discrets.

\section{Notions générales}

\subsection{Classification}

Considérons une équation aux dérivées partielles d'ordre 2 de la forme
\[
  A\frac{\partial^2 U}{\partial x^2} + 2B\frac{\partial^2 U}{\partial x\,\partial y}
  + C\frac{\partial^2 U}{\partial y^2}
  + F\Bigl(x, y, U, \frac{\partial U}{\partial x}, \frac{\partial U}{\partial y}\Bigr) = 0
\]
où $A$, $B$ et $C$ sont des fonctions de $x$ et $y$ et sont de classe $C^2$.
\begin{itemize}[label=$\blacklozenge$]
  \item Si $B^2 - AC = 0 \;\Longrightarrow\;$ équation \emph{parabolique} ;
  \item Si $B^2 - AC < 0 \;\Longrightarrow\;$ équation \emph{elliptique} ;
  \item Si $B^2 - AC > 0 \;\Longrightarrow\;$ équation \emph{hyperbolique}.
\end{itemize}

Pour résoudre une équation aux dérivées partielles, il faut connaître les
conditions aux limites ou aux frontières et les conditions initiales. La
nature de ces conditions permet de définir les types de problèmes suivants :

\textbf{Problème de DIRICHLET pur :} les valeurs de $U$ sur les frontières du
domaine sont connues.

\textbf{Problème de NEUMANN pur :} les valeurs des dérivées de $U$ sur les
frontières sont connues.

\textbf{Problème mixte DIRICHLET-NEUMANN :} les valeurs de $U$ sont connues
sur une partie des frontières et celles des dérivées de $U$ sur la partie
restante.

\textbf{Problème de CAUCHY :} les valeurs de $U$ et de ses dérivées sont
connues sur les frontières du domaine.

\subsection{Quelques équations aux dérivées partielles en Physique}

\subsubsection{Équation d'onde}

Elle a la forme :
\[
  \frac{\partial^2 U}{\partial t^2} = a^2\Delta U
  + f\Bigl(\frac{\partial U}{\partial t}, x, y, z, t\Bigr)
\]
où $a$ est la vitesse de l'onde et $f$ une action extérieure.

\textbf{En dimension 1}, cette équation décrit :
\begin{itemize}[label=--]
  \item Les vibrations transversales d'une corde : l'équation générale dans ce
        cas est
        \[
          \rho(x)\frac{\partial^2 U}{\partial t^2} + k\frac{\partial U}{\partial t}
          = T\frac{\partial^2 U}{\partial x^2} + f(x,t)
        \]
        où $\rho$ est la masse volumique, $k$ le coefficient de dissipation,
        $T$ la tension dans la corde. Ici la vitesse des ondes est
        $a = \sqrt{T/\rho}$.
  \item Les vibrations longitudinales dans une barre : $a = \sqrt{E/\rho}$ où
        $E$ est le module d'élasticité.
  \item Les vibrations de torsion d'une barre avec $a = \sqrt{\mu/\rho}$ où
        $\mu = \dfrac{E}{2(1+\sigma)}$ et $\sigma$ est le coefficient de
        Poisson.
  \item La propagation d'ondes électromagnétiques ou électriques.
\end{itemize}

\textbf{En dimension 2}, elle peut décrire les vibrations d'une membrane
plane.

\textbf{En dimension 3}, elle décrit les ondes sonores en hydrodynamique, les
ondes électromagnétiques (cas des cavités résonantes et du rayonnement des
antennes).

\subsubsection{Équation de Laplace et de Poisson}

Elles s'écrivent sous la forme suivante
\[
  \begin{aligned}
    \Delta U &= 0 && \text{(Laplace)}\\
    \Delta U &= f(x, y, z) && \text{(Poisson)}
  \end{aligned}
\]
Ce sont les prototypes des équations elliptiques. Elles décrivent les
problèmes stationnaires suivants :
\begin{itemize}[label=--]
  \item répartition des charges électriques sur la surface d'un conducteur ;
  \item équilibre thermique dans un corps homogène ;
  \item distribution du potentiel des vitesses dans un fluide non visqueux,
        incompressible et en écoulement irrotationnel.
\end{itemize}

\subsubsection{Équation de la chaleur}

Sa forme générale est :
\[
  c\rho\frac{\partial U}{\partial t} = \operatorname{div}\bigl(k\,\overrightarrow{\operatorname{grad}}\,U\bigr)
  + F(x, y, z, t)
\]
où $U$ est la température en un point $P(x, y, z)$ au temps $t$, $\rho$ la
masse volumique, $c$ la chaleur spécifique, $k$ la conductivité interne et $F$
l'apport des sources ou des puits de chaleur.

Si $c$, $\rho$ et $k$ sont constantes, cette équation se réduit à :
\[
  \frac{\partial U}{\partial t} = a^2\Delta U + f(x, y, z, t) .
\]
Dans le cas où
\[
  \frac{\partial U}{\partial t} = 0,
\]
on est en présence d'un problème de transfert de chaleur en régime permanent
(conduction et convection).

Dans le cas général où $\dfrac{\partial U}{\partial t} \neq 0$, on est en
présence d'un problème de diffusion de la chaleur.

\section{Équation de Poisson dans le plan}

\subsection{Discrétisation de l'équation de Poisson}

Soit l'équation :
\begin{equation}
  \frac{\partial^2 U}{\partial x^2} + \frac{\partial^2 U}{\partial y^2} = f(x, y)
  \tag{P}
\end{equation}
définie sur un domaine $D$ rectangulaire :
$D = \{(x, y) \in \R^2 \,/\, a < x < b \text{ et } c < y < d\}$.

Soit $(S)$ la frontière de $D$, la condition suivante :
\[
  U(x, y) = g(x, y)
\]
est imposée sur $(S)$.

Nous considérons que les fonctions $f$ et $g$ sont continues sur le domaine
$D$. Soient $h$ et $k$ les pas de calcul suivant $x$ et $y$ respectivement.
Soient $n$ et $m$ les nombres entiers définis par
\[
  h = \frac{b - a}{n} \qquad\text{et}\qquad k = \frac{d - c}{m} .
\]
Un point du domaine $D$ est défini par les coordonnées suivantes :
\[
  \left\{
  \begin{aligned}
    x_i &= a + ih && \text{avec } i = 1, \dots, n-1\\
    y_j &= c + jk && \text{avec } j = 1, \dots, m-1
  \end{aligned}
  \right.
\]
D'après les définitions des dérivées, on a :
\begin{equation}
  \frac{\partial^2 U}{\partial x^2}
  = \frac{U_{i+1,j} - 2U_{i,j} + U_{i-1,j}}{h^2}
  - \frac{h^2}{12}\frac{\partial^4 U}{\partial x^4}(\xi_{ij}, y_j)
  \tag{P1}
\end{equation}
où $U_{ij} = U(x_i, y_j)$ ; $\xi \in\, ]x_{i-1}, x_{i+1}[$.

En écrivant une équation similaire pour $\dfrac{\partial^2 U}{\partial y^2}$,
l'équation (P) devient, avec une erreur de troncature d'ordre
$O(h^2 + k^2)$, le système suivant :
\begin{equation}
  2\left[\Bigl(\frac{h}{k}\Bigr)^2 + 1\right]U_{i,j} - \bigl(U_{i+1,j} + U_{i-1,j}\bigr)
  - \Bigl(\frac{h}{k}\Bigr)^2\bigl(U_{i,j+1} + U_{i,j-1}\bigr) = -h^2f_{i,j}
  \tag{P2}
\end{equation}
où $f_{i,j} = f(x_i, y_j)$ ; avec $i = 1, 2, \dots, n-1$ et
$j = 1, 2, \dots, m-1$.

La discrétisation des conditions aux frontières est la suivante :
\begin{equation}
  \left\{
  \begin{aligned}
    U_{0,j} &= g(x_0, y_j)\\
    U_{n,j} &= g(x_n, y_j)
  \end{aligned}
  \right.
  \quad j = 0, 1, \dots, m
  \qquad\qquad
  \left\{
  \begin{aligned}
    U_{i,0} &= g(x_i, y_0)\\
    U_{i,m} &= g(x_i, y_m)
  \end{aligned}
  \right.
  \quad i = 0, 1, \dots, n
  \tag{P3}
\end{equation}
En associant les équations (P2) et (P3), on obtient le système d'équations
algébriques linéaires en $U_{ij}$ que l'on peut résoudre par la méthode
d'élimination de Gauss.

\begin{exercice}
Écrire le schéma d'itération de Gauss-Seidel et de Jacobi pour résoudre
l'équation de Poisson.
\end{exercice}

\subsection{Cas spécifiques}

\subsubsection{Conditions aux limites de type Neumann}

Lorsque les conditions aux limites sont de type Neumann, le schéma de
discrétisation doit se faire avec beaucoup de précautions. Généralement, il
est conseillé de combiner cette condition avec l'équation aux dérivées
partielles de départ.

Par exemple, si la condition de Neumann s'écrit
$\left.\dfrac{\partial U}{\partial y}\right|_S = g(x, y)$, alors on utilise le
schéma suivant. On développe $U(x_i, y_1)$ autour du point $(x_i, y_0)$ ;
c'est-à-dire que l'on écrit :
\[
  U(x_i, y_1) = U(x_i, y_0) + k\frac{\partial U}{\partial y}(x_i, y_0)
  + \frac{k^2}{2}\frac{\partial^2 U}{\partial y^2}(x_i, y_0) + O(k^3) .
\]
La quantité $\dfrac{\partial^2 U}{\partial y^2}$ en $(x_i, y_0)$ peut être
remplacée par son équivalent à partir de l'équation aux dérivées partielles
de départ. Par exemple, dans le cas de l'équation de Laplace, on a :
$\dfrac{\partial^2 U}{\partial y^2} = -\dfrac{\partial^2 U}{\partial x^2}$.
Il vient alors :
\[
  \begin{aligned}
    U(x_i, y_1) &= U(x_i, y_0) + k\frac{\partial U}{\partial y}(x_i, y_0)
                 - \frac{k^2}{2}\frac{\partial^2 U}{\partial x^2}(x_i, y_0)\\
                &= U(x_i, y_0) + k\frac{\partial U}{\partial y}(x_i, y_0)
                 - \frac{k^2}{2h^2}\bigl(U_{i+1,0} - 2U_{i,0} + U_{i-1,0}\bigr).
  \end{aligned}
\]
Finalement la condition de Neumann donne :
\begin{equation}
  U_{i,1} - \Bigl(1 + \frac{k^2}{h^2}\Bigr)U_{i,0}
  + \frac{k^2}{2h^2}\bigl(U_{i+1,0} + U_{i-1,0}\bigr) = kg(x_i, y_0) .
  \tag{P4}
\end{equation}
On peut alors associer à (P4) l'équation (P2) pour avoir le système discret
global à résoudre.

\subsubsection{Théorème du maximum et du minimum discrets}

Dans le cas de l'équation de Laplace, on a :
\[
  U_{i,j} = \frac{k^2}{2(k^2 + h^2)}\bigl(U_{i+1,j} + U_{i-1,j}\bigr)
          + \frac{h^2}{2(k^2 + h^2)}\bigl(U_{i,j+1} + U_{i,j-1}\bigr).
\]
Donc $U_{ij}$ est le barycentre des nombres $U_{i+1,j}$, $U_{i-1,j}$,
$U_{i,j+1}$ et $U_{i,j-1}$ affectés des poids respectifs :
$\alpha = \dfrac{k^2}{2(k^2 + h^2)}$, $\alpha$,
$\beta = \dfrac{h^2}{2(k^2 + h^2)}$ et $\beta$. Puisque ces poids sont tous
positifs, $U_{ij}$ est donc compris entre la plus grande et la plus petite des
quantités $U_{i+1,j}$, $U_{i-1,j}$, $U_{i,j+1}$ et $U_{i,j-1}$. D'où le
théorème suivant :

\paragraph{\underline{Théorème}} Dans le cas de l'équation de Laplace, il n'y
a pas de maximum et de minimum stricts de $U$ à l'intérieur du domaine $D$. Le
maximum et le minimum sont atteints sur les frontières de $D$.

\paragraph{\underline{Remarque}} D'autres schémas de différences finies
peuvent être établis à partir des formules de discrétisation des dérivées. Par
exemple, on peut définir $\dfrac{\partial^2 U}{\partial x^2}$ par des schémas
présentant des erreurs d'ordre $O(h^4 + k^4)$.

\section{Équation de la chaleur}

\subsection{Problème à une dimension d'espace}

Nous nous limiterons au cas unidimensionnel. On a l'équation :
\begin{equation}
  \frac{\partial U}{\partial t} = a^2\frac{\partial^2 U}{\partial x^2}
  \tag{C1}
\end{equation}
dans le domaine $D = T \times L$ ; où $T = \R_+$ et $L = \,]0, l[$.

Les conditions aux limites sont : $U(0, t) = 0$, $U(l, t) = 0$.

Et la condition initiale est
\[
  U(x, 0) = f(x) .
\]
Il s'agit par exemple d'une tige de longueur $l$ dont les extrémités $x = 0$
et $x = l$ sont maintenues à la température $U = 0$. À l'instant initial, on
la soumet à une distribution de température $U(x, 0) = f(x)$. Il est question
de suivre l'évolution temporelle de la température dans la poutre.

\subsection{Méthode des différences progressives ou schéma explicite}

Soit $h$ le pas spatial et $k$ le pas temporel et soit $m$ le nombre défini
par $mh = l$.

On a
\[
  \frac{\partial U}{\partial t} = \frac{U_{i,j+1} - U_{i,j}}{k}
  \qquad\text{et}\qquad
  \frac{\partial^2 U}{\partial x^2} = \frac{U_{i+1,j} - 2U_{i,j} + U_{i-1,j}}{h^2} .
\]
Dans le domaine $D$, on obtient alors le système discret suivant :
\begin{equation}
  U_{i,j+1} = \Bigl(1 - \frac{2a^2k}{h^2}\Bigr)U_{i,j}
            + \frac{a^2k}{h^2}\bigl(U_{i+1,j} + U_{i-1,j}\bigr) ;
  \quad i = 1, 2, \dots, m-1 ; \quad j = 0, 1, 2, \dots
  \tag{C2}
\end{equation}
Le schéma (C2) présente une erreur de troncature d'ordre $O(k + h^2)$.

C'est un schéma explicite car pour calculer $U_{i,j+1}$, il suffit de
remplacer les quantités figurant au second membre par leurs valeurs connues et
stockées en mémoire.

Les conditions initiales et aux limites se discrétisent de la manière
suivante :
\begin{equation}
  U_{i,0} = f(x_i), \; i = 0, 1, \dots, m \,; \qquad
  U_{0,j} = 0 \;\text{ et }\; U_{m,j} = 0, \; j = 0, 1, 2, \dots
  \tag{C3}
\end{equation}
La première condition (C3) peut être substituée dans (C2) pour déterminer
$U_{i,1}$ pour $i$ allant de $1$ jusqu'à $m-1$. La condition $U_{0,j} = 0$ et
$U_{m,j} = 0 \Longrightarrow U_{0,1} = U_{m,1} = 0$. On a alors toutes les
valeurs de $U_{i,1}$. En procédant de la même manière, on obtient les valeurs
de $U_{i,2}$, $U_{i,3}$ et ainsi de suite.

L'inconvénient de cette méthode est qu'elle est conditionnellement stable. Il
faut donc avoir une contrainte entre les pas $h$ et $k$ et les coefficients de
l'équation pour obtenir des résultats satisfaisants. C'est la condition de
stabilité que nous établissons dans le paragraphe qui suit.

\subsection{Condition de stabilité du schéma explicite}

Considérons le schéma aux différences finies progressives de l'équation de la
chaleur. Soit une perturbation $\varepsilon_{ij}$ de la solution numérique
$U_{i,j}$. La solution numérique devient alors
$U'_{i,j} = U_{i,j} + \varepsilon_{ij}$. En écrivant que $U'_{ij}$ vérifie le
schéma discret, on obtient que la perturbation $\varepsilon_{ij}$ vérifie
aussi l'équation discrète :
\[
  \frac{\varepsilon_{i,j+1} - \varepsilon_{i,j}}{k}
  = a^2\frac{\varepsilon_{i+1,j} - 2\varepsilon_{i,j} + \varepsilon_{i-1,j}}{h^2} .
\]
Le schéma discret est stable si $\varepsilon_{ij}$ reste borné ou tend vers 0
au fur et à mesure que le temps avance. Pour avoir la condition de stabilité,
nous supposons que
\[
  \varepsilon_{i,j} = \sin(iph)\,e^{-rjk} .
\]
Cette expression est utilisée en tenant compte de sa dépendance suivant
l'espace $x = ih$ et doit satisfaire aux conditions aux limites. En
introduisant cette expression dans l'équation vérifiée par $\varepsilon_{ij}$,
on obtient :
\[
  e^{-rk} = 1 - \frac{4ka^2}{h^2}\sin^2\frac{ph}{2} .
\]
Or $\varepsilon_{i,j+1} = \varepsilon_{i,j}\,e^{-rk}$.

Donc $e^{-rk}$ est le \emph{facteur d'amplification}.

Le schéma est stable si le facteur d'amplification est inférieur à 1 (en
module). Dans le cas particulier analysé ici, cette condition implique que :
\[
  \abs{1 - \frac{4ka^2}{h^2}\sin^2\frac{ph}{2}} < 1
  \;\Longrightarrow\;
  \abs{1 - \frac{4ka^2}{h^2}} < 1
  \;\Longrightarrow\;
  \frac{ka^2}{h^2} < \frac{1}{2} .
\]

\subsection{Autres schémas}

\subsubsection{Différences finies régressives ou schéma implicite}

Afin d'éviter l'utilisation d'une condition de stabilité, on peut utiliser les
différences finies régressives pour la dérivée temporelle ; c'est-à-dire
\begin{equation}
  \frac{\partial U}{\partial t} = \frac{U_{i,j} - U_{i,j-1}}{k} .
  \tag{C4}
\end{equation}
L'équation de la chaleur (C1) donne alors le système discret suivant :
\begin{equation}
  (1 + 2\lambda)\,U_{i,j} - \lambda U_{i+1,j} - \lambda U_{i-1,j} = U_{i,j-1}
  \qquad\text{pour}\quad 1 \le i \le m-1 \quad\text{et}\quad 1 \le j \le \infty
  \tag{C5}
\end{equation}
où $\lambda = a^2k/h^2$.

Compte tenu des conditions initiales et des conditions aux limites, on obtient
un système matriciel tri-diagonal que l'on peut résoudre par les méthodes
itératives ou la méthode d'élimination de Gauss.

Pour déterminer la condition de stabilité, on procède de la même façon que
précédemment. On trouve pour ce schéma implicite, que le facteur
d'amplification est
\[
  e^{-rk} = \frac{1}{1 + \dfrac{4ka^2}{h^2}\sin^2\dfrac{ph}{2}} .
\]
On note qu'il est toujours inférieur à 1. Par conséquent, ce schéma est
inconditionnellement stable.

\subsubsection{Méthode de Richardson}

Les méthodes ci-dessus ont une erreur de troncature d'ordre $O(k)$ en $k$. Si
on veut obtenir une erreur d'ordre $O(k^2)$, on utilise la formule suivante :
\begin{equation}
  \frac{\partial U}{\partial t} = \frac{U_{i,j+1} - U_{i,j-1}}{2k} .
  \tag{C6}
\end{equation}
Avec cette discrétisation de la dérivée temporelle, on aboutit à un schéma
instable. Il n'est donc pas utilisé dans le cadre de l'équation de la
chaleur.

\subsubsection{Méthode de Crank-Nicolson}

Cette méthode fait la somme du schéma des différences finies progressives et
du schéma des différences régressives. Mais il faut au préalable écrire la
formule des différences régressives à l'instant $j+1$. On obtient alors
\[
  \frac{U_{i,j+1} - U_{i,j}}{k}
  - \frac{a^2}{2h^2}\bigl(U_{i+1,j} - 2U_{i,j} + U_{i-1,j}
                         + U_{i+1,j+1} - 2U_{i,j+1} + U_{i-1,j+1}\bigr) = 0 .
\]
C'est une méthode plus précise qui présente une erreur d'ordre
$O(h^2 + k^2)$.

\subsubsection{Méthode de Dufort-Frankel}

Pour résoudre le problème d'instabilité de la méthode de Richardson, la
quantité $U_{ij}$ est remplacée par $\bigl(U_{i,j+1} + U_{i,j-1}\bigr)/2$ et
on obtient
\[
  \frac{U_{i,j+1} - U_{i,j-1}}{2k}
  - a^2\bigl(U_{i+1,j} - U_{i,j+1} - U_{i,j-1} + U_{i-1,j}\bigr)/h^2 = 0 .
\]
Ce schéma est inconditionnellement stable.

\begin{exercice}[Exercices]
\begin{enumerate}[label=\arabic*)]
  \item Établir l'algorithme pour le schéma explicite.
  \item Établir les conditions de stabilité des schémas de Crank-Nicolson et
        de Dufort-Frankel.
  \item Établir les algorithmes pour les méthodes des différences
        régressives, de Crank-Nicolson et de Dufort-Frankel.
\end{enumerate}
\end{exercice}

\section{Équation d'ondes}

Nous considérons le problème suivant :
\begin{equation}
  \frac{\partial^2 U}{\partial t^2} - a^2\frac{\partial^2 U}{\partial x^2} = 0
  \tag{O1}
\end{equation}
dans le domaine $D = T \times L$ et soumis aux conditions suivantes :
\[
  U(0, t) = U(l, t) = 0 \,; \qquad U(x, 0) = f(x) \,; \qquad
  \frac{\partial U(x, 0)}{\partial t} = g(x) .
\]
Il s'agit par exemple d'une corde vibrante fixée à ses deux extrémités et
excitée en une région par un déplacement transversal $U(x, 0) = f(x)$ et une
vitesse de déplacement transversal $\dfrac{\partial U(x, 0)}{\partial t} = g(x)$.
La question est de savoir comment cette perturbation va se propager le long de
la corde.

\subsection{Discrétisation directe}

La discrétisation de (O1) par les différences finies donne :
\begin{equation}
  U_{i,j+1} = 2(1 - \lambda^2)\,U_{i,j} + \lambda^2\bigl(U_{i+1,j} + U_{i-1,j}\bigr) - U_{i,j-1}
  \tag{O2}
\end{equation}
où $\lambda = ak/h$, $m = l/h$ et $i = 1, \dots, m-1$ ; $j = 1, \dots, \infty$.

Les conditions aux limites donnent :
\begin{equation}
  U_{0,j} = U_{m,j} = 0 .
  \tag{O3}
\end{equation}
La première condition initiale donne
\begin{equation}
  U_{i,0} = f(x_i) .
  \tag{O4}
\end{equation}
La deuxième condition initiale permet d'écrire :
\begin{equation}
  U_{i,1} = U_{i,0} + k\,g(x_i) .
  \tag{O5}
\end{equation}
Mais (O5) présente une erreur de troncature d'ordre $O(k)$. Pour avoir une
erreur d'ordre $O(k^2)$ comme l'équation (O2), on procède de la manière
suivante.
\[
  \frac{\partial U(x_i, 0)}{\partial t}
  = \frac{U(x_i, t_1) - U(x_i, 0)}{k}
  - \frac{k}{2}\frac{\partial^2 U(x_i, 0)}{\partial t^2} + O(k^2) .
\]
Or
\[
  \frac{\partial^2 U(x_i, 0)}{\partial t^2}
  = a^2\frac{\partial^2 U(x_i, 0)}{\partial x^2}
  = a^2\frac{\partial^2 f}{\partial x^2}
  = a^2\frac{f_{i+1} - 2f_i + f_{i-1}}{h^2} .
\]
En utilisant la condition initiale, on obtient finalement :
\begin{equation}
  U_{i,1} = (1 - \lambda^2)\,f_i + \frac{\lambda^2}{2}\bigl(f_{i+1} + f_{i-1}\bigr) + kg_i
  \tag{O6}
\end{equation}
avec $i = 1, \dots, m-1$.

Il s'agit donc de résoudre le système discret formé des équations (O2) et
(O6). La recherche des valeurs $U_{i,j+1}$ nécessite la connaissance des
valeurs $U_{i,j}$ et $U_{i,j-1}$. Au début, il faut connaître $U_{i,0}$ et
$U_{i,1}$. Les valeurs de $U_{i,0}$ découlent de (O4) et celles de $U_{i,1}$
découlent de (O6). On peut ainsi lancer le processus itératif pour calculer
toutes les valeurs de $U_{ij}$.

Il faut noter que ce schéma de discrétisation directe a une condition de
stabilité qui est
\[
  \lambda < 1 \,; \quad\text{c'est-à-dire}\quad ak/h < 1 .
\]

\begin{exercice}
\begin{enumerate}[label=\alph*-]
  \item Établir la condition de stabilité du schéma explicite ci-dessus.
  \item Établir l'algorithme du schéma explicite ci-dessus.
\end{enumerate}
\end{exercice}

\subsection{Décomposition en un système d'équations du premier ordre}

Dans certains cas, il est avantageux de décomposer l'équation d'ondes en deux
équations du premier ordre avant de discrétiser.

Par exemple, en posant
\[
  V = \frac{\partial U}{\partial x} \qquad\text{et}\qquad
  W = \frac{1}{a}\frac{\partial U}{\partial t},
\]
l'équation d'ondes devient le système d'équations suivant
\[
  \left\{
  \begin{aligned}
    \frac{\partial V}{\partial t} &= a\frac{\partial W}{\partial x}\\[0.4em]
    \frac{\partial W}{\partial t} &= a\frac{\partial V}{\partial x}
  \end{aligned}
  \right.
\]
Ce système d'équations peut être discrétisé de manière explicite et de manière
implicite.

\subsubsection{Discrétisation explicite}

Pour la discrétisation explicite, on peut utiliser la discrétisation par les
différences finies progressives dans le temps et par les différences centrées
dans l'espace. On obtient alors
\[
  \left\{
  \begin{aligned}
    V_{i,j+1} &= V_{i,j} + \frac{ak}{2h}\bigl(W_{i+1,j} - W_{i-1,j}\bigr)\\[0.3em]
    W_{i,j+1} &= W_{i,j} + \frac{ak}{2h}\bigl(V_{i+1,j} - V_{i-1,j}\bigr)
  \end{aligned}
  \right.
\]
Pour établir la condition de stabilité, posons
\[
  \left\{
  \begin{aligned}
    V_{i,j} &= A_j\,e^{Ipih}\\
    W_{i,j} &= B_j\,e^{Ipih}
  \end{aligned}
  \right.
  \qquad\text{où } I^2 = -1 .
\]
En substituant dans le système discret, on obtient
\[
  \left\{
  \begin{aligned}
    V_{i,j+1} &= \Bigl[A_j + B_j\,\frac{Iak}{h}\sin ph\Bigr]e^{Ipih}\\[0.3em]
    W_{i,j+1} &= \Bigl[B_j + A_j\,\frac{Iak}{h}\sin ph\Bigr]e^{Ipih}
  \end{aligned}
  \right.
\]
or
\[
  \left\{
  \begin{aligned}
    V_{i,j+1} &= A_{j+1}\,e^{Ipih}\\
    W_{i,j+1} &= B_{j+1}\,e^{Ipih}
  \end{aligned}
  \right.
\]
En comparant ces expressions, on obtient
\[
  \begin{pmatrix} A_{j+1}\\ B_{j+1} \end{pmatrix}
  = \begin{pmatrix} 1 & Ic\\ Ic & 1 \end{pmatrix}
    \begin{pmatrix} A_{j}\\ B_{j} \end{pmatrix}
  \qquad\text{où}\quad c = \frac{ak}{h}\sin ph .
\]
La matrice d'amplification est donc
\[
  A = \begin{pmatrix} 1 & Ic\\ Ic & 1 \end{pmatrix}.
\]
La stabilité du schéma exige que les valeurs propres de la matrice
d'amplification $A$ aient des modules inférieurs à 1. Dans le cas présent, les
modules des valeurs propres de $A$ sont tous supérieurs à 1. Donc le schéma
est inconditionnellement instable.

\subsubsection{Amélioration du schéma explicite}

Compte tenu de cette instabilité inconditionnelle, on peut améliorer le schéma
naturel ci-dessus. On peut trouver tous les $V_{i,j+1}$ de la première
équation et les utiliser ensuite dans la deuxième équation. On obtient alors le
schéma discret suivant :
\[
  \left\{
  \begin{aligned}
    V_{i,j+1} &= V_{i,j} + \frac{ak}{2h}\bigl(W_{i+1,j} - W_{i-1,j}\bigr)\\[0.3em]
    W_{i,j+1} &= W_{i,j} + \frac{ak}{2h}\bigl(V_{i+1,j+1} - V_{i-1,j+1}\bigr)
  \end{aligned}
  \right.
\]
Pour un tel schéma, la matrice d'amplification est
\[
  A = \begin{pmatrix} 1 & Ic\\ Ic & 1 - c^2 \end{pmatrix}.
\]
On prouve que le schéma est stable si et seulement si $\dfrac{ak}{h} < 2$.

\subsubsection{Schéma implicite}

On peut avoir un schéma implicite en discrétisant les dérivées spatiales à
l'instant $j+1$ et les dérivées temporelles par les différences progressives.
On obtient alors
\[
  \left\{
  \begin{aligned}
    V_{i,j+1} &= V_{i,j} + \frac{ak}{2h}\bigl(W_{i+1,j+1} - W_{i-1,j+1}\bigr)\\[0.3em]
    W_{i,j+1} &= W_{i,j} + \frac{ak}{2h}\bigl(V_{i+1,j+1} - V_{i-1,j+1}\bigr)
  \end{aligned}
  \right.
\]
Les valeurs propres de la matrice d'amplification sont dans ce cas
$1/(1 + Ic)$, $1/(1 - Ic)$.

Elles sont toutes de module $< 1$. Donc ce schéma implicite est
inconditionnellement stable.

\subsubsection{Schéma de Lax}

Si dans le schéma explicite naturel (premier schéma), on remplace $V_{i,j}$ et
$W_{i,j}$ par leurs valeurs moyennes respectives entre $i-1$ et $i+1$, alors
on obtient le schéma de Lax qui s'écrit
\[
  \left\{
  \begin{aligned}
    V_{i,j+1} &= \frac{1}{2}\bigl(V_{i+1,j} + V_{i-1,j}\bigr)
               + \frac{ak}{2h}\bigl(W_{i+1,j} - W_{i-1,j}\bigr)\\[0.3em]
    W_{i,j+1} &= \frac{1}{2}\bigl(W_{i+1,j} + W_{i-1,j}\bigr)
               + \frac{ak}{2h}\bigl(V_{i+1,j} - V_{i-1,j}\bigr)
  \end{aligned}
  \right.
\]
Dans ce cas la matrice d'amplification est
\[
  A = \begin{pmatrix} \cos ph & Ic\\ Ic & \cos ph \end{pmatrix}.
\]
La condition de stabilité est $ak/h < 1$.

\subsubsection{Schéma de Lax-Wendroff}

Pour ce schéma, on développe $V_{i,j+1}$ et $W_{i,j+1}$ jusqu'à l'ordre 2.
Par exemple
\[
  V_{i,j+1} = V_{i,j} + k\frac{\partial V_{i,j}}{\partial t}
            + \frac{k^2}{2}\frac{\partial^2 V_{i,j}}{\partial t^2}
            = V_{i,j} + ka\frac{\partial W_{i,j}}{\partial x}
            + \frac{k^2a^2}{2}\frac{\partial^2 V_{i,j}}{\partial x^2} .
\]
En faisant la même chose sur $W$, on obtient le schéma discret de
Lax-Wendroff suivant :
\[
  \left\{
  \begin{aligned}
    V_{i,j+1} &= V_{i,j} + \frac{ak}{2h}\bigl(W_{i+1,j} - W_{i-1,j}\bigr)
               + \frac{a^2k^2}{2h^2}\bigl(V_{i+1,j} - 2V_{i,j} + V_{i-1,j}\bigr)\\[0.3em]
    W_{i,j+1} &= W_{i,j} + \frac{ak}{2h}\bigl(V_{i+1,j} - V_{i-1,j}\bigr)
               + \frac{a^2k^2}{2h^2}\bigl(W_{i+1,j} - 2W_{i,j} + W_{i-1,j}\bigr)
  \end{aligned}
  \right.
\]
La matrice d'amplification est
\[
  A = \begin{pmatrix}
    1 + \dfrac{a^2k^2}{h^2}(\cos ph - 1) & \dfrac{Iak}{h}\sin ph\\[1em]
    \dfrac{Iak}{h}\sin ph & 1 + \dfrac{a^2k^2}{h^2}(\cos ph - 1)
  \end{pmatrix}.
\]
La condition de stabilité : $\dfrac{ak}{h} < 1$.

\paragraph{Remarque} Il faut noter cependant que lorsqu'une équation
hyperbolique présente des conditions initiales discontinues, la méthode des
différences finies n'est plus adaptée ; on doit faire appel à d'autres
méthodes telle que la méthode des caractéristiques.

\chapter*{Exercices}
\addcontentsline{toc}{chapter}{Exercices}
\markboth{Exercices}{}

\subsection*{Exercice 1}

La répartition de l'énergie rayonnée par une étoile en fonction de la
longueur d'onde est donnée par la loi de Planck
\[
  I(\lambda) = (8\pi hc)\Big/\left(\lambda^5\left[\exp\Bigl(\frac{hc}{\lambda kT}\Bigr) - 1\right]\right)
\]
où $h$, $k$ et $c$ sont respectivement les constantes de Planck, de Boltzmann
et la célérité de la lumière. On pose $x = hc/\lambda kT$. D'après les
résultats expérimentaux, $I$ est maximale pour $x_m$.
\begin{enumerate}[label=\arabic*-]
  \item Établir l'équation vérifiée par $x_m$.
  \item On veut trouver $x_m$ par la méthode de Newton-Raphson. Établir
        l'algorithme permettant de calculer la valeur numérique de $x_m$.
\end{enumerate}

\subsection*{Exercice 2}

L'équation de Van der Waals pour une mole de gaz de masse $M$ est
\[
  \Bigl(P + \frac{a}{V^2}\Bigr)(V - b) = RMT
\]
où $a$, $b$ et $R$ sont des constantes. $P$, $T$ et $V$ sont respectivement la
pression, la température et le volume.

Établir un algorithme qui permet de calculer le volume $V$ chaque fois que
l'on connaît la température et la pression.

\subsection*{Exercice 3}

En 1976, Desantis a établi que le facteur de compressibilité $b$ des gaz réels
vérifie la relation
\[
  z = \bigl(1 + y + y^2 - y^3\bigr)\big/\bigl(1 - y^3\bigr)
\]
où $y = b/4v$ ; $v$ est le volume molaire et $z$ est une constante. On veut
trouver $b$ par la méthode de Newton-Raphson.
\begin{enumerate}[label=\arabic*-]
  \item Expliquer le principe de la méthode.
  \item Établir alors un algorithme permettant de résoudre le problème.
  \item Traduire cet algorithme en Fortran.
\end{enumerate}

\subsection*{Exercice 4}

Une grandeur $y$ vérifie l'équation suivante :
\[
  \frac{\dd y}{\dd x} = -ay - by^4 + c
  \qquad\text{avec}\quad y(0) = y_0 .
\]
$a$, $b$ et $c$ sont des constantes.

On se propose de résoudre cette équation numériquement par la méthode de
prédicteur-correcteur (Adams-Bashforth du second ordre associé à la méthode
de Runge-Kutta d'ordre 2).
\begin{enumerate}[label=\arabic*)-]
  \item Établir le schéma itératif correspondant.
  \item Établir l'algorithme de résolution de l'équation.
  \item Traduire cet algorithme en Pascal et en Fortran.
\end{enumerate}

\subsection*{Exercice 5}

Dans le domaine $x \in\, ]0, 1[$ une grandeur $y$ vérifie l'équation
différentielle
\[
  y'' + \frac{a}{x}y' + be^{cy} = 0
  \qquad\text{avec}\quad y'(0) = 0 \quad\text{et}\quad y(1) = 1 .
\]
$a$, $b$ et $c$ sont des constantes. On veut résoudre cette équation par la
méthode de tir (associée à la méthode de Runge-Kutta d'ordre 4).
\begin{enumerate}[label=\arabic*)-]
  \item En utilisant la méthode de Newton (pour les équations algébriques non
        linéaires), donner le schéma itératif qui permet de calculer les
        valeurs successives de la condition de tir en $x = 0$.
  \item Établir la méthode itérative de Runge-Kutta d'ordre 4.
  \item Établir l'algorithme complet (incluant le schéma des conditions
        initiales) permettant de résoudre le problème.
\end{enumerate}

\subsection*{Exercice 6}

Dans l'espace $KA$ entre une cathode et une anode $A$ parallèle à $K$, le
potentiel à la distance $x$ de $K$ satisfait l'équation différentielle
\begin{equation}
  V'' = aV^{-1/2}
  \tag{1}
\end{equation}
où $a$ est une constante. On se propose de trouver numériquement les valeurs
de $V$.
\begin{enumerate}[label=\arabic*-]
  \item Dans cette première question, on suppose que le potentiel et le champ
        électrique sont nuls sur la cathode (problème à valeurs initiales). On
        résout alors (1) par la méthode de Runge-Kutta d'ordre 2.
        \begin{enumerate}[label=\alph*-]
          \item Sachant que RK2 dérive de la formule intégrale des trapèzes,
                donner l'expression générale des itérations RK2 pour une
                équation différentielle du premier ordre.
          \item L'équation (1) peut se décomposer en deux équations
                différentielles du premier ordre. Établir l'algorithme de
                résolution numérique de (1) par RK2.
          \item Traduire cet algorithme en Pascal et en Fortran.
        \end{enumerate}
  \item Dans cette deuxième question, on suppose connues les conditions aux
        limites $V_K = 0$ et $V_A = V_0$, la distance $KA = l$. Le problème
        peut alors se résoudre par la méthode de tir.
        \begin{enumerate}[label=\alph*-]
          \item Expliquer le principe de la méthode de tir.
          \item Écrire alors l'algorithme de résolution de (1) par la méthode
                de tir.
        \end{enumerate}
\end{enumerate}

\subsection*{Exercice 7}

Soit une équation de la forme :
\[
  y''' = f(x, y, y', y'')
  \qquad\text{avec}\quad y(a) = y_1 \quad\text{et}\quad y(b) = y_2 .
\]
On veut résoudre cette équation par la méthode de tir.
\begin{enumerate}[label=\arabic*)-]
  \item Établir le problème aux valeurs initiales résultant.
  \item Le problème aux valeurs initiales résultant doit être résolu par la
        méthode prédicteur-correcteur.
        \begin{enumerate}[label=\alph*)-]
          \item Expliquer le principe de la méthode prédicteur-correcteur sur
                une équation différentielle du premier ordre.
          \item En déduire l'algorithme de résolution du problème
                différentiel du troisième ordre ci-dessus.
        \end{enumerate}
\end{enumerate}

\subsection*{Exercice 8}

\begin{enumerate}[label=\arabic*-]
  \item On considère une fonction $y$ de classe $C^4$ sur un intervalle
        $[a,b]$.
        \begin{enumerate}[label=\alph*-]
          \item En utilisant la formule d'interpolation polynomiale de degré 2,
                établir que $y' = \bigl(y(x+h) - y(x-h)\bigr)/2h$ où $h$ est
                le pas de dérivation.
          \item Déduire de a) une expression pour $y''$.
        \end{enumerate}
  \item Soit le problème aux valeurs aux limites
        $y'' = p(x)y' + q(x)y + g(x)$ défini sur $[a,b]$ avec $y(a) = c$ et
        $y(b) = d$. On veut résoudre ce problème par la méthode des
        différences finies.

        Établir à partir de la question 1) le système algébrique découlant du
        problème.
  \item On veut résoudre le système algébrique ci-dessus par la méthode
        itérative de Jacobi.
        \begin{enumerate}[label=\alph*-]
          \item Expliquer le principe de la méthode itérative de Jacobi.
          \item Écrire alors l'algorithme permettant de résoudre le problème.
          \item Traduire l'algorithme en Pascal.
        \end{enumerate}
\end{enumerate}

\subsection*{Exercice 9}

Une particule de masse $M$ se déplace dans un potentiel de la forme :
\[
  u(y) = \frac{1}{2}ky^2 + \frac{1}{4}k_1y^4 .
\]
Le milieu présente une viscosité linéaire de coefficient $\alpha$ et la
particule est soumise à une force extérieure $f(t) = F\cos\omega t$.
\begin{enumerate}[label=\arabic*)-]
  \item Établir l'équation du mouvement de la particule et adimensionner cette
        équation (rendre l'équation sans dimension).
  \item On donne $y(0) = \dot{y}(0) = 0$ et on veut utiliser la méthode de
        Runge-Kutta d'ordre 4 pour trouver les valeurs de $y$ et $\dot{y}$ en
        fonction du temps $t$.

        Établir l'algorithme de résolution numérique de l'équation
        adimensionnée trouvée à la question (1).
  \item À partir de la résolution numérique, on obtient un ensemble de valeurs
        discrètes $y_i = y(t_i)$, l'expression analytique $y(t)$ étant
        inconnue, on propose la formule approchée suivante :
        \[
          y(t) = \sum_{n=1}^{N}\bigl(a_n\cos n\omega t + b_n\sin n\omega t\bigr).
        \]
        Expliquer comment on peut utiliser la méthode des moindres carrés pour
        déterminer les valeurs des coefficients $a_n$ et $b_n$. Établir les
        algorithmes permettant de déterminer ces coefficients.
\end{enumerate}

\subsection*{Exercice 10}

On considère un système algébrique linéaire $n \times n$ de la forme $AX = b$
où $A = (a_{ij})$, $b = (b_i)$ et $X = (x_i)$ avec $i, j = 1, 2, \dots, n$.

On veut résoudre ce système par la méthode d'élimination de Gauss ou du
pivot.
\begin{enumerate}[label=\arabic*-]
  \item Indiquer le principe de la méthode.
  \item Établir l'algorithme correspondant.
  \item De l'algorithme ci-dessus, on obtient un système triangulaire. Établir
        l'algorithme de résolution du système triangulaire obtenu.
  \item Traduire ce dernier algorithme en Pascal.
\end{enumerate}

\subsection*{Exercice 11}

Le mouvement d'un pendule est décrit par l'équation différentielle
\[
  \frac{\dd^2 Y}{\dd t^2} + \omega_0^2\sin Y = 0
\]
où $\omega_0$ est la pulsation propre du pendule et $Y$ le déplacement
angulaire.
\begin{enumerate}[label=\arabic*-]
  \item On veut résoudre numériquement cette équation par la méthode de
        Runge-Kutta d'ordre 2.
        \begin{enumerate}[label=\alph*-]
          \item Établir l'algorithme de résolution de cette équation.
          \item Traduire cet algorithme en Fortran.
        \end{enumerate}
  \item Lorsque le pendule oscille entre $Y_0$ et $-Y_0$, sa période est
        définie par la relation
        \begin{equation}
          T(Y_0) = \frac{2T_0}{\pi}\int_0^{\pi/2}\frac{\dd x}{\sqrt{1 - A(\sin x)^2}}
          \tag{2}
        \end{equation}
        où $T_0 = 2\pi/\omega_0$ et $A = \sin^2(Y_0/2)$.
        \begin{enumerate}[label=\alph*-]
          \item Après avoir établi la formule composite des trapèzes, donner
                l'algorithme de calcul de l'intégrale (2).
          \item Traduire cet algorithme en Pascal et en Fortran.
        \end{enumerate}
  \item De l'intégration numérique de l'intégrale (2), on obtient le tableau
        suivant :
        \begin{center}
          \begin{tabular}{|l|c|c|c|c|}
            \hline
            $Y_0$ & $\pi/18$ & $\pi/9$ & $\pi/6$ & $2\pi/9$\\ \hline
            $T(Y_0)/T_0$ & 1,002 & 1,008 & 1,020 & 1,030\\ \hline
          \end{tabular}
        \end{center}
        On exprime alors $T(Y_0)$ sous la forme $T(Y_0) = T_0(a + bY_0^2)$. En
        utilisant la méthode des moindres carrés, déterminer les valeurs
        numériques de $a$ et de $b$.
\end{enumerate}

\subsection*{Exercice 12}

On mesure la capacité calorifique $C_v$ d'un corps à volume constant et on
obtient les résultats suivants
\begin{center}
  \begin{tabular}{lccc}
    $T$ (K) & 0.1 & 0.3 & 0.5\\
    $C_v \times 10^4$ cal/K.mole & 0.211 & 0.694 & 1.361
  \end{tabular}
\end{center}
\begin{enumerate}[label=\arabic*-]
  \item On se propose de déterminer la chaleur $Q$ absorbée par une mole du
        corps lorsque sa température passe de 0.1~K à 0.5~K ; calculer $Q$ par
        la formule de Simpson.
  \item En réalité, le tableau ci-dessus est une partie d'un tableau plus
        général contenant de nombreuses valeurs expérimentales. On propose la
        formule $C_v = a + bT$. On veut alors déterminer les coefficients $a$
        et $b$ par la méthode des moindres carrés.
        \begin{enumerate}[label=\alph*-]
          \item Définir le système d'équations vérifié par $a$ et $b$.
          \item Établir des algorithmes permettant d'évaluer numériquement les
                coefficients du système d'équations et de calculer les
                valeurs de $a$ et de $b$.
          \item Traduire ces algorithmes en Pascal.
          \item On veut calculer $Q$ par la formule composite de Simpson à
                partir des valeurs expérimentales de $C_v$. Établir
                l'algorithme de calcul.
        \end{enumerate}
\end{enumerate}

\subsection*{Exercice 13}

On considère l'équation intégrale :
\begin{equation}
  \phi(x) = f(x) + \int_a^b k(x,t)\,\phi(t)\,\dd t
  \tag{1}
\end{equation}
où $k(x,t)$ est une fonction connue. On subdivise le domaine $[a,b]$ en $n$
points $t_i = ih$ équidistants.
\begin{enumerate}[label=\arabic*)-]
  \item L'équation (1) peut être approximée par la formule composite des
        trapèzes. Elle prend alors la forme :
        \begin{equation}
          \phi(x) = f(x) + (b - a)\sum_{i=1}^{n} c_i\,k(x, t_i)\,\phi_i
          \tag{2}
        \end{equation}
        où $\phi_i = \phi(t_i)$.

        Donner les expressions des coefficients $c_i$.
  \item Sachant que l'équation (2) est vérifiée en tout point $x_j = t_j$,
        montrer que l'équation intégrale (1) se réduit au système
        algébrique :
        \begin{equation}
          \phi_j = f_j + (b - a)\sum_{i=1}^{n} c_i\,k(t_j, t_i)\,\phi_i
          \tag{3}
        \end{equation}
        où les $\phi_j$ et $\phi_i$ sont les inconnues.
  \item Le système algébrique (3) peut se mettre sous la forme :
        \begin{equation}
          AX = b
          \tag{4}
        \end{equation}
        où $A$ est une matrice $(n,n)$, $X$ et $b$ sont les matrices colonnes
        de $n$ éléments. Définir ces matrices.
  \item On veut résoudre le système (4) par la méthode de Gauss.
        \begin{enumerate}[label=\alph*)-]
          \item Donner l'algorithme de triangularisation du système.
          \item Écrire l'algorithme de résolution du système triangularisé.
          \item Traduire ce dernier algorithme en Pascal et Fortran.
        \end{enumerate}
  \item $\phi(x)$ possède un zéro dans l'intervalle $[a,b]$ et on se propose
        de trouver ce zéro en utilisant la méthode de bissection. Décrire
        cette méthode et donner son algorithme.
\end{enumerate}

\subsection*{Exercice 14}

À l'instant initial $t = 0$, la température d'un mur d'épaisseur $L$ vaut
$T_0$. Sur la surface $x = 0$, la température varie suivant la loi :
\[
  T(x = 0, t) = T_0 + T_1\sin\frac{\pi}{2}t \qquad\text{pour } t > 0 .
\]
Par un mécanisme approprié, la face $x = L$ est maintenue à la température
constante $T_L$. L'équation de la chaleur dans le milieu est alors :
\[
  \frac{\partial T}{\partial t} = \alpha\frac{\partial^2 T}{\partial x^2}
  \qquad\text{où } \alpha \text{ est la diffusibilité.}
\]
\begin{enumerate}[label=\arabic*)-]
  \item On se propose de résoudre numériquement le problème ainsi posé par la
        méthode des différences finies progressives.
        \begin{enumerate}[label=\alph*)-]
          \item Donner la forme discrète du problème. Quel est l'ordre de
                l'erreur de discrétisation ?
          \item Écrire l'algorithme de résolution du système discret obtenu.
          \item Traduire cet algorithme en Pascal.
        \end{enumerate}
  \item On veut maintenant utiliser la méthode de Richardson pour résoudre le
        même problème.
        \begin{enumerate}[label=\alph*)-]
          \item Donner le nouveau système discret et indiquer l'ordre et
                l'erreur de discrétisation.
          \item En admettant que l'on utilise le schéma discret de la question
                (1) pour approximer les températures
                $T_{i,1} = T(x = i\Delta_x, \Delta t)$, établir l'algorithme
                permettant de trouver la température aux instants ultérieurs.
          \item Traduire cet algorithme en Fortran.
        \end{enumerate}
\end{enumerate}

\subsection*{Exercice 15}

\begin{enumerate}[label=\arabic*)]
  \item On considère l'équation de la chaleur
        $\dfrac{\partial u}{\partial t} = a^2\dfrac{\partial^2 u}{\partial x^2}$
        définie sur $D = \R_+ \times\, ]0, l[$ avec
        $u(0,t) = u(l,t) = 0$, $u(x,0) = f(x)$.

        On veut résoudre cette équation par le schéma explicite avec
        $\dfrac{\partial u}{\partial t} = \dfrac{u_{i,j+1} - u_{i,j}}{k}$.

        Établir que la condition de stabilité de ce schéma est
        $\dfrac{a^2k}{h^2} \le \dfrac{1}{2}$ où $k$ et $h$ sont les pas
        temporel et spatial.
  \item Que devient cette condition de stabilité dans le cas d'un problème
        plan : $\dfrac{\partial u}{\partial t} = a^2\Delta u$ ? On prendra
        $h_1$ et $h_2$ les pas spatiaux.
  \item On considère l'équation d'advection-diffusion dans le plan :
        \[
          \frac{\partial\Omega}{\partial t} + u_1\frac{\partial\Omega}{\partial x}
          + u_2\frac{\partial\Omega}{\partial y} = \nu\Delta\Omega .
        \]
        On la discrétise par le schéma centré en espace. Établir que la
        condition de stabilité est $\dfrac{\nu k}{h^2} \le \dfrac{1}{4}$ où
        $k$ est le pas temporel et $h$ est le pas spatial suivant $x$ et sur
        $y$.

        Que devient cette condition dans le cas unidimensionnel ?
\end{enumerate}

\subsection*{Exercice 16}

Soit l'équation d'onde suivante :
\[
  \frac{\partial^2 u}{\partial t^2} + a\frac{\partial u}{\partial t}
  = c^2\frac{\partial^2 u}{\partial x^2} + du^3
\]
où $(t, x) \in \R_+ \times\, ]0, l[$ et $u(0,t) = u(l,t) = 0$,
$u(x,0) = f(x)$, $\dfrac{\partial u}{\partial t}(x,0) = g(x)$.

$a$, $c$ et $d$ sont des constantes.
\begin{enumerate}[label=\arabic*)-]
  \item Définir chaque terme de cette équation.
  \item On veut résoudre numériquement cette équation par les méthodes de
        différences finies centrées. Donner la forme discrète du problème
        ainsi posé.
  \item Exprimer $u_{i,2} = u(x = ih, t = 2k)$ en fonction des valeurs des
        fonctions $f$ et $g$ aux points $i-2$ ; $i-1$ ; $i$ ; $i+1$ et $i+2$
        et des pas $h$ et $k$.
  \item Établir un algorithme de résolution du système discret obtenu à la
        question (2).
  \item Traduire l'algorithme dans un langage de programmation de votre choix.
\end{enumerate}

\subsection*{Exercice 17}

\begin{center}
  \begin{tikzpicture}[>=Stealth]
    % coupe du tube (1 cm pour 125 mm) : région hachurée entre les deux rectangles
    \fill[pattern=north east lines, even odd rule]
      (0,0) rectangle (8,6.4) (2,1.6) rectangle (6,4.8);
    \draw[thick] (0,0) rectangle (8,6.4);
    \draw[thick] (2,1.6) rectangle (6,4.8);
    \node at (3.2,3.6) {$p_s$};
    % cotes intérieures
    \draw[<->, dashed] (2,2.4) -- (6,2.4) node[pos=0.5, below] {$L_1$};
    \draw[<->, dashed] (5.2,1.6) -- (5.2,4.8) node[pos=0.6, left] {$l_1$};
    % cotes extérieures
    \draw[<->] (-0.5,0) -- (-0.5,6.4) node[midway, left] {$l$};
    \draw[<->] (0,-0.5) -- (8,-0.5) node[midway, below] {$L$};
  \end{tikzpicture}
\end{center}
\[
  L = 1000~\text{mm}, \quad L_1 = 500~\text{mm}, \quad l = 800~\text{mm}, \quad l_1 = 400~\text{mm}
\]
Le dispositif représente la coupe d'un tube rectangulaire dans lequel
s'écoule dans sa partie centrale un liquide avec une pression $p_s$. La
pression sur la surface extérieure est nulle. À chaque point de la région
hachurée, la distribution des pressions vérifie l'équation
\[
  \frac{\partial^2 p}{\partial x^2} + \frac{\partial^2 p}{\partial y^2} = 0 .
\]
\begin{enumerate}[label=\arabic*)-]
  \item Faites un adimensionnement de cette équation et précisez les nouvelles
        conditions sur la frontière. La pression adimensionnée est $u$.
  \item \begin{enumerate}[label=\alph*)-]
          \item En utilisant la méthode des différences finies, établir
                l'équation discrète vérifiée par $u_{i,j}$ en tout point
                $(x_i = ih, y_j = jk)$.
          \item Discrétiser les conditions aux limites.
        \end{enumerate}
  \item On se propose de trouver les $u_{i,j}$ par la méthode itérative de
        Gauss-Seidel.
        \begin{enumerate}[label=\alph*)-]
          \item Décrire cette méthode itérative.
          \item Établir l'algorithme de calcul des $u_{i,j}$.
        \end{enumerate}
  \item On suppose que le tube a maintenant une forme cylindrique de rayon
        intérieur $R_1 = 400$~mm et de rayon extérieur $R_2 = 800$~mm.
        \begin{enumerate}[label=\alph*)-]
          \item Quelle est alors l'équation vérifiée par la pression $p$ ou la
                grandeur adimensionnée $u$ ?
          \item Proposer une méthode de discrétisation de cette nouvelle
                équation.
        \end{enumerate}
\end{enumerate}

\subsection*{Exercice 18}

\begin{center}
  \begin{tikzpicture}[scale=0.6,>=Stealth]
    % plaque L x l avec une encoche centrale de largeur L/3 et de hauteur l/2
    \fill[pattern=north east lines]
      (0,0) -- (3,0) -- (3,3) -- (6,3) -- (6,0) -- (9,0) -- (9,6) -- (0,6) -- cycle;
    \draw[thick]
      (0,0) -- (3,0) -- (3,3) -- (6,3) -- (6,0) -- (9,0) -- (9,6) -- (0,6) -- cycle;
    % cotes
    \draw[<->] (0,6.6) -- (9,6.6) node[midway, above] {$L$};
    \draw[<->] (-0.6,0) -- (-0.6,6) node[midway, left] {$l$};
    \draw[<->] (0,-0.6) -- (3,-0.6) node[midway, below] {$\frac{L}{3}$};
    \draw[<->] (6,-0.6) -- (9,-0.6) node[midway, below] {$\frac{L}{3}$};
    \draw[<->] (5.5,0) -- (5.5,3) node[midway, left] {$\frac{l}{2}$};
  \end{tikzpicture}
\end{center}

La répartition d'une grandeur physique $u$ sur la plaque représentée vérifie
l'équation de Poisson suivante
\[
  \frac{\partial^2 u}{\partial x^2} + \frac{\partial^2 u}{\partial y^2} = f(x, y)
\]
avec $u = g(x, y)$ sur la frontière.
\begin{enumerate}[label=\arabic*)-]
  \item En utilisant la méthode des différences finies, établir l'équation
        discrète vérifiée par $u_{i,j}$.
  \item On se propose de trouver les $u_{i,j}$ par la méthode itérative de
        Jacobi.
        \begin{enumerate}[label=\alph*)-]
          \item Décrire cette méthode.
          \item Établir l'algorithme de calcul de $u_{i,j}$.
        \end{enumerate}
\end{enumerate}

\subsection*{Exercice 19}

\begin{center}
  \begin{tikzpicture}[scale=0.8,>=Stealth]
    \fill[pattern=north east lines] (0,0) -- (4,0) -- (0,4) -- cycle;
    \draw[thick] (0,0) -- (4,0) -- (0,4) -- cycle;
    \draw[<->] (-0.5,0) -- (-0.5,4) node[midway, left] {$h$};
    \draw[<->] (0,-0.5) -- (4,-0.5) node[midway, below] {$l$};
  \end{tikzpicture}
\end{center}

La répartition d'une grandeur physique sur la plaque triangulaire ci-dessus
vérifie l'équation de Laplace :
$\dfrac{\partial^2 u}{\partial x^2} + \dfrac{\partial^2 u}{\partial y^2} = 0$
avec $u = 0$ sur l'hypoténuse et $\dfrac{\partial u}{\partial n} = a$ sur les
autres côtés ($n$ est la normale par rapport aux côtés considérés).
\begin{enumerate}[label=\arabic*)-]
  \item Discrétiser ce problème par la méthode des différences finies.
  \item Établir l'algorithme de calcul des valeurs discrètes $u_{i,j}$ sur
        toute la plaque.
\end{enumerate}


\end{document}
